\documentclass[10pt,compress,t,notes=noshow,xcolor=table]{beamer}

% =====================================================================
%  Interpretable Machine Learning: Evaluation of iML
%  One 90-minute session in 7 chunks. Timing estimates per chunk are in
%  the comment banners above each \section. Total: approx. 89 min.
%
%  Style handling
%  * Inside the course repository (../../style/preamble.tex exists), the
%    course style and the latex-math macros are loaded unchanged.
%  * Otherwise, a minimal standalone fallback defines the same commands
%    (framev, \splitVTT, \titlemeta, itemizeM, \spacer, ...) so that this
%    file compiles on its own.
%  Session-specific macros (citations, colors, placeholders) are defined
%  after that block and are used in both modes.
% =====================================================================

\newif\ifcoursestyle
\IfFileExists{../../style/preamble.tex}{\coursestyletrue}{\coursestylefalse}

\ifcoursestyle
  % graphicx and color are loaded via lmu-lecture.sty
% maxwidth is the original width if it is less than linewidth
% otherwise use linewidth (to make sure the graphics do not exceed the margin)
% TODO: Remove once cleared to be superfluous
% \makeatletter
% \def\maxwidth{ %
%   \ifdim\Gin@nat@width>\linewidth
%     \linewidth
%   \else
%     \Gin@nat@width
%   \fi
% }
% \makeatother

% ---------------------------------%
% latex-math dependencies, do not remove:
% - mathtools
% - bm
% - siunitx
% - dsfont
% - xspace
% ---------------------------------%

%--------------------------------------------------------%
%       Language, encoding, typography
%--------------------------------------------------------%

\usepackage[english]{babel}
\usepackage[utf8]{inputenc} % Enables inputting UTF-8 symbols
% Standard AMS suite (loaded via lmu-lecture.sty)

% Font for double-stroke / blackboard letters for sets of numbers (N, R, ...)
% Distribution name is "doublestroke"
% According to https://mirror.physik.tu-berlin.de/pub/CTAN/fonts/doublestroke/dsdoc.pdf
% the "bbm" package does a similar thing and may be superfluous.
% Required for latex-math
\usepackage{dsfont}

% bbm – "Blackboard-style" cm fonts (https://www.ctan.org/pkg/bbm)
% Used to be in common.tex, loaded directly after this file
% Maybe superfluous given dsfont is loaded
% TODO: Check if really unused?
% \usepackage{bbm}

% bm – Access bold symbols in maths mode - https://ctan.org/pkg/bm
% Required for latex-math, preferred over \boldsymbol
% https://tex.stackexchange.com/questions/3238/bm-package-versus-boldsymbol
\usepackage{bm}

% pifont – Access to PostScript standard Symbol and Dingbats fonts
% Used for \newcommand{\xmark}{\ding{55}, which is never used
% aside from lecture_advml/attic/xx-automl/slides.Rnw
% \usepackage{pifont}

% Quotes (inline and display), provdes \enquote
% https://ctan.org/pkg/csquotes
\usepackage{csquotes}

% Adds arg to enumerate env, technically superseded by enumitem according
% to https://ctan.org/pkg/enumerate
% Replace with https://ctan.org/pkg/enumitem ?
% Even better: enumitem is not really compatible with beamer and breaks all sorts of things
% particularly the enumerate environment. The enumerate package also just isn't required
% from what I can tell so... don't re-add it I guess?
% \usepackage{enumerate}

% Line spacing - provides \singlespacing \doublespacing \onehalfspacing
% https://ctan.org/pkg/setspace
% \usepackage{setspace}

% mathtools – Mathematical tools to use with amsmath
% https://ctan.org/pkg/mathtools?lang=en
% latex-math dependency according to latex-math repo
\usepackage{mathtools}

% Maybe not great to use this https://tex.stackexchange.com/a/197/19093
% Use align instead -- TODO: Global search & replace to check, eqnarray is used a lot
% $ rg -f -u "\begin{eqnarray" -l | grep -v attic | awk -F '/' '{print $1}' | sort | uniq -c
%   13 lecture_advml
%   14 lecture_i2ml
%    2 lecture_iml
%   27 lecture_optimization
%   45 lecture_sl
\usepackage{eqnarray}
% For shaded regions / boxes
% Used sometimes in optim
% https://www.ctan.org/pkg/framed
\usepackage{framed}

%--------------------------------------------------------%
%       Cite button (version 2024-05)
%--------------------------------------------------------%

% Superseded by style/ref-buttons.sty, kept just in case these don't work out somehow.

% Note this requires biber to be in $PATH when running,
% telltale error in log would be e.g. Package biblatex Info: ... file 'authoryear.dbx' not found
% aside from obvious "biber: command not found" or similar.
% Tried moving this to lmu-lecture.sty but had issues I didn't quite understood,
% so it's here for now.

\usepackage{hyperref}

% Only try adding a references file if it exists, otherwise
% this would compile error when references.bib is not found
% NOTE: Bibliography packages (usebib, biblatex) are now loaded by ref-buttons.sty when needed
% This keeps all bibliography-related setup in one place

% Legacy \citelink command removed - superseded by ref-buttons.sty

%--------------------------------------------------------%
%       Displaying code and algorithms
%--------------------------------------------------------%

% Reimplements verbatim environments: https://ctan.org/pkg/verbatim
% verbatim used sed at least once in
% supervised-classification/slides-classification-tasks.tex
% Removed since code should not be put on slides anyway
% \usepackage{verbatim}

% Both used together for algorithm typesetting, see also overleaf: https://www.overleaf.com/learn/latex/Algorithms
% algorithmic env is also used, but part of the bundle:
%   "algpseudocode is part of the algorithmicx bundle, it gives you an improved version of algorithmic besides providing some other features"
% According to https://tex.stackexchange.com/questions/229355/algorithm-algorithmic-algorithmicx-algorithm2e-algpseudocode-confused
\usepackage{algorithm}
\usepackage{algpseudocode}

%--------------------------------------------------------%
%       Tables
%--------------------------------------------------------%

% multi-row table cells: https://www.namsu.de/Extra/pakete/Multirow.html
% Provides \multirow
% Used e.g. in evaluation/slides-evaluation-measures-classification.tex
\usepackage{multirow}

% colortbl: https://ctan.org/pkg/colortbl
% "The package allows rows and columns to be coloured, and even individual cells." well.
% Provides \columncolor and \rowcolor
% \rowcolor is used multiple times, e.g. in knn/slides-knn.tex
\usepackage{colortbl}

% long/multi-page tables: https://texdoc.org/serve/longtable.pdf/0
% Not used in slides
% \usepackage{longtable}

% pretty table env: https://ctan.org/pkg/booktabs
% Is used
% Defines \toprule
\usepackage{booktabs}

%--------------------------------------------------------%
%       Figures: Creating, placing, verbing
%--------------------------------------------------------%

% wrapfig - Wrapping text around figures https://de.overleaf.com/learn/latex/Wrapping_text_around_figures
% Provides wrapfigure environment -used in lecture_optimization
\usepackage{wrapfig}

% Sub figures in figures and tables
% https://ctan.org/pkg/subfig -- supersedes subfigure package
% Provides \subfigure
% \subfigure not used in slides but slides-tuning-practical.pdf errors without this pkg, error due to \captionsetup undefined
\usepackage{subfig}

% Actually it's pronounced PGF https://en.wikibooks.org/wiki/LaTeX/PGF/TikZ
\usepackage{tikz}

% No idea what/why these settings are what they are but I assume they're there on purpose
\usetikzlibrary{shapes,arrows,automata,positioning,calc,chains,trees, shadows}
\tikzset{
  %Define standard arrow tip
  >=stealth',
  %Define style for boxes
  punkt/.style={
    rectangle,
    rounded corners,
    draw=black, very thick,
    text width=6.5em,
    minimum height=2em,
    text centered},
  % Define arrow style
  pil/.style={
    ->,
    thick,
    shorten <=2pt,
    shorten >=2pt,}
}

%--------------------------------------------------------%
%       Beamer setup and custom macros & environments
%--------------------------------------------------------%

% Main sty file for beamer setup (layout, style, lecture page numbering, etc.)
% For long-term maintenance, this may me refactored into a more modular set of .sty files
\usepackage{../../style/lmu-lecture}
% Custom itemize wrappers, itemizeS, itemizeL, etc
\usepackage{../../style/customitemize}
% Custom framei environment, uses custom itemize!
\usepackage{../../style/framei}
% Custom frame2 environment, allows specifying font size for all content
\usepackage{../../style/frame2}
% Column layout macros
\usepackage{../../style/splitV}
% \image and derivatives
\usepackage{../../style/image}
% New generation of reference button macros
\usepackage{../../style/ref-buttons}

% Used regularly
\let\code=\texttt

% Not sure what/why this does
\setkeys{Gin}{width=0.9\textwidth}

% -- knitr leftovers --
% These may be used by knitr/R Markdown workflows in other lectures
\makeatletter
\def\maxwidth{ %
  \ifdim\Gin@nat@width>\linewidth
    \linewidth
  \else
    \Gin@nat@width
  \fi
}
\makeatother

% Define colors for syntax highlighting (may be used by knitr)
\definecolor{fgcolor}{rgb}{0.345, 0.345, 0.345}
\definecolor{shadecolor}{rgb}{.97, .97, .97}

% knitr code output environment
\newenvironment{knitrout}{}{}


% Can't find a reason why common.tex is not just part of this file?
\usepackage[]{graphicx}
\usepackage[]{color}
% maxwidth is the original width if it is less than linewidth
% otherwise use linewidth (to make sure the graphics do not exceed the margin)
\makeatletter
\def\maxwidth{ %
\ifdim\Gin@nat@width>\linewidth
\linewidth
\else
\Gin@nat@width
\fi
}
\makeatother
%\usepackage[fontsize=10.5pt]{scrextend}
\definecolor{ggred}{rgb}{0.973, 0.463, 0.427}
\definecolor{ggblue}{rgb}{0, 0.749, 0.769}
\definecolor{fgcolor}{rgb}{0.345, 0.345, 0.345}
\newcommand{\hlnum}[1]{\textcolor[rgb]{0.686,0.059,0.569}{#1}}%
\newcommand{\hlstr}[1]{\textcolor[rgb]{0.192,0.494,0.8}{#1}}%
\newcommand{\hlcom}[1]{\textcolor[rgb]{0.678,0.584,0.686}{\textit{#1}}}%
\newcommand{\hlopt}[1]{\textcolor[rgb]{0,0,0}{#1}}%
\newcommand{\hlstd}[1]{\textcolor[rgb]{0.345,0.345,0.345}{#1}}%
\newcommand{\hlkwa}[1]{\textcolor[rgb]{0.161,0.373,0.58}{\textbf{#1}}}%
\newcommand{\hlkwb}[1]{\textcolor[rgb]{0.69,0.353,0.396}{#1}}%
\newcommand{\hlkwc}[1]{\textcolor[rgb]{0.333,0.667,0.333}{#1}}%
\newcommand{\hlkwd}[1]{\textcolor[rgb]{0.737,0.353,0.396}{\textbf{#1}}}%
\newcommand{\predvar}{Var\left[\hat{f}(\xv)\right]}
\let\hlipl\hlkwb

\usepackage{pdfpages}
\usepackage{framed}
\makeatletter
\newenvironment{kframe}{%
\def\at@end@of@kframe{}%
\ifinner\ifhmode%
\def\at@end@of@kframe{\end{minipage}}%
\begin{minipage}{\columnwidth}%
\fi\fi%
\def\FrameCommand##1{\hskip\@totalleftmargin \hskip-\fboxsep
\colorbox{shadecolor}{##1}\hskip-\fboxsep
% There is no \\@totalrightmargin, so:
\hskip-\linewidth \hskip-\@totalleftmargin \hskip\columnwidth}%
\MakeFramed {\advance\hsize-\width
\@totalleftmargin\z@ \linewidth\hsize
\@setminipage}}%
{\par\unskip\endMakeFramed%
\at@end@of@kframe}
\makeatother

\definecolor{shadecolor}{rgb}{.97, .97, .97}
\definecolor{messagecolor}{rgb}{0, 0, 0}
\definecolor{warningcolor}{rgb}{1, 0, 1}
\definecolor{errorcolor}{rgb}{1, 0, 0}
\newenvironment{knitrout}{}{} % an empty environment to be redefined in TeX

\usepackage{alltt}
\newcommand{\SweaveOpts}[1]{}  % do not interfere with LaTeX
\newcommand{\SweaveInput}[1]{} % because they are not real TeX commands
\newcommand{\Sexpr}[1]{}       % will only be parsed by R

\usepackage[english]{babel}
\usepackage[utf8]{inputenc}

\usepackage{dsfont}
\usepackage{verbatim}
\usepackage{amsmath}
\usepackage{amsfonts}
\usepackage{bm}
\usepackage{csquotes}
\usepackage{multirow}
\usepackage{longtable}
\usepackage{booktabs}
\usepackage{enumerate}
\usepackage[absolute,overlay]{textpos}
\usepackage{psfrag}
\usepackage{algorithm}
\usepackage{algpseudocode}
\usepackage{eqnarray}
\usepackage{arydshln}
\usepackage{tabularx}
\usepackage{placeins}
\usepackage{tikz}
\usepackage{setspace}
\usepackage{colortbl}
\usepackage{mathtools}
\usepackage{wrapfig}
\usepackage{bm}
\usepackage[backend=biber]{biblatex}

\usetikzlibrary{tikzmark, shapes,arrows,automata,positioning,calc,chains,trees,  shadows, decorations.pathreplacing}
\tikzset{
%Define standard arrow tip
>=stealth',
%Define style for boxes
punkt/.style={
rectangle,
rounded corners,
draw=black, very thick,
text width=6.5em,
minimum height=2em,
text centered},
% Define arrow style
pil/.style={
->,
thick,
shorten <=2pt,
shorten >=2pt,}
}

\usepackage{subfig}

\usepackage{bbm}
% basic latex stuff
\newcommand{\pkg}[1]{{\fontseries{b}\selectfont #1}} %fontstyle for R packages
\newcommand{\lz}{\vspace{0.5cm}} %vertical space
\newcommand{\dlz}{\vspace{1cm}} %double vertical space
\newcommand{\oneliner}[1] % Oneliner for important statements
{\begin{block}{}\begin{center}\begin{Large}#1\end{Large}\end{center}\end{block}}

% Latexmath Notation
% math spaces
\ifdefined\N                                                                
\renewcommand{\N}{\mathds{N}} % N, naturals
\else \newcommand{\N}{\mathds{N}} \fi 
\newcommand{\Z}{\mathds{Z}} % Z, integers
\newcommand{\Q}{\mathds{Q}} % Q, rationals
\newcommand{\R}{\mathds{R}} % R, reals
\ifdefined\C 
  \renewcommand{\C}{\mathds{C}} % C, complex
\else \newcommand{\C}{\mathds{C}} \fi
\newcommand{\continuous}{\mathcal{C}} % C, space of continuous functions
\newcommand{\M}{\mathcal{M}} % machine numbers
\newcommand{\epsm}{\epsilon_m} % maximum error

% counting / finite sets
\newcommand{\setzo}{\{0, 1\}} % set 0, 1
\newcommand{\setmp}{\{-1, +1\}} % set -1, 1
\newcommand{\unitint}{[0, 1]} % unit interval

% basic math stuff
\newcommand{\xt}{\tilde x} % x tilde
\newcommand{\argmax}{\operatorname{arg\,max}} % argmax
\newcommand{\argmin}{\operatorname{arg\,min}} % argmin
\newcommand{\argminlim}{\mathop{\mathrm{arg\,min}}\limits} % argmax with limits
\newcommand{\argmaxlim}{\mathop{\mathrm{arg\,max}}\limits} % argmin with limits  
\newcommand{\sign}{\operatorname{sign}} % sign, signum
\newcommand{\I}{\mathbb{I}} % I, indicator
\newcommand{\order}{\mathcal{O}} % O, order
\newcommand{\pd}[2]{\frac{\partial{#1}}{\partial #2}} % partial derivative
\newcommand{\floorlr}[1]{\left\lfloor #1 \right\rfloor} % floor
\newcommand{\ceillr}[1]{\left\lceil #1 \right\rceil} % ceiling

% sums and products
\newcommand{\sumin}{\sum\limits_{i=1}^n} % summation from i=1 to n
\newcommand{\sumim}{\sum\limits_{i=1}^m} % summation from i=1 to m
\newcommand{\sumjn}{\sum\limits_{j=1}^n} % summation from j=1 to p
\newcommand{\sumjp}{\sum\limits_{j=1}^p} % summation from j=1 to p
\newcommand{\sumik}{\sum\limits_{i=1}^k} % summation from i=1 to k
\newcommand{\sumkg}{\sum\limits_{k=1}^g} % summation from k=1 to g
\newcommand{\sumjg}{\sum\limits_{j=1}^g} % summation from j=1 to g
\newcommand{\meanin}{\frac{1}{n} \sum\limits_{i=1}^n} % mean from i=1 to n
\newcommand{\meankg}{\frac{1}{g} \sum\limits_{k=1}^g} % mean from k=1 to g
\newcommand{\prodin}{\prod\limits_{i=1}^n} % product from i=1 to n
\newcommand{\prodkg}{\prod\limits_{k=1}^g} % product from k=1 to g
\newcommand{\prodjp}{\prod\limits_{j=1}^p} % product from j=1 to p

% linear algebra
\newcommand{\one}{\boldsymbol{1}} % 1, unitvector
\newcommand{\zero}{\mathbf{0}} % 0-vector
\newcommand{\id}{\boldsymbol{I}} % I, identity
\newcommand{\diag}{\operatorname{diag}} % diag, diagonal
\newcommand{\trace}{\operatorname{tr}} % tr, trace
\newcommand{\spn}{\operatorname{span}} % span
\newcommand{\scp}[2]{\left\langle #1, #2 \right\rangle} % <.,.>, scalarproduct
\newcommand{\mat}[1]{\begin{pmatrix} #1 \end{pmatrix}} % short pmatrix command
\newcommand{\Amat}{\mathbf{A}} % matrix A
\newcommand{\Deltab}{\mathbf{\Delta}} % error term for vectors

% basic probability + stats
\renewcommand{\P}{\mathds{P}} % P, probability
\newcommand{\E}{\mathds{E}} % E, expectation
\newcommand{\var}{\mathsf{Var}} % Var, variance
\newcommand{\cov}{\mathsf{Cov}} % Cov, covariance
\newcommand{\corr}{\mathsf{Corr}} % Corr, correlation
\newcommand{\normal}{\mathcal{N}} % N of the normal distribution
\newcommand{\iid}{\overset{i.i.d}{\sim}} % dist with i.i.d superscript
\newcommand{\distas}[1]{\overset{#1}{\sim}} % ... is distributed as ...
% machine learning
\newcommand{\Xspace}{\mathcal{X}} % X, input space
\newcommand{\Yspace}{\mathcal{Y}} % Y, output space
\newcommand{\Zspace}{\mathcal{Z}} % Z, space of sampled datapoints
\newcommand{\nset}{\{1, \ldots, n\}} % set from 1 to n
\newcommand{\pset}{\{1, \ldots, p\}} % set from 1 to p
\newcommand{\gset}{\{1, \ldots, g\}} % set from 1 to g
\newcommand{\Pxy}{\mathbb{P}_{xy}} % P_xy
\newcommand{\Exy}{\mathbb{E}_{xy}} % E_xy: Expectation over random variables xy
\newcommand{\xv}{\mathbf{x}} % vector x (bold)
\newcommand{\Xv}{\mathbf{X}} % vector X (random variable, bold)
\newcommand{\xtil}{\tilde{\mathbf{x}}} % vector x-tilde (bold)
\newcommand{\yv}{\mathbf{y}} % vector y (bold)
\newcommand{\Yv}{\mathbf{Y}} % vector Y (random variable, bold)
\newcommand{\xy}{(\xv, y)} % observation (x, y)
\newcommand{\xvec}{\left(x_1, \ldots, x_p\right)^\top} % (x1, ..., xp)
\newcommand{\Xmat}{\mathbf{X}} % Design matrix
\newcommand{\allDatasets}{\mathds{D}} % The set of all datasets
\newcommand{\allDatasetsn}{\mathds{D}_n}  % The set of all datasets of size n
\newcommand{\D}{\mathcal{D}} % D, data
\newcommand{\Dn}{\D_n} % D_n, data of size n
\newcommand{\Dtrain}{\mathcal{D}_{\text{train}}} % D_train, training set
\newcommand{\Dtest}{\mathcal{D}_{\text{test}}} % D_test, test set
\newcommand{\xyi}[1][i]{\left(\xv^{(#1)}, y^{(#1)}\right)} % (x^i, y^i), i-th observation
\newcommand{\Dset}{\left( \xyi[1], \ldots, \xyi[n]\right)} % {(x1,y1)), ..., (xn,yn)}, data
\newcommand{\defAllDatasetsn}{(\Xspace \times \Yspace)^n} % Def. of the set of all datasets of size n
\newcommand{\defAllDatasets}{\bigcup_{n \in \N}(\Xspace \times \Yspace)^n} % Def. of the set of all datasets
\newcommand{\xdat}{\left\{ \xv^{(1)}, \ldots, \xv^{(n)}\right\}} % {x1, ..., xn}, input data
\newcommand{\ydat}{\left\{ \yv^{(1)}, \ldots, \yv^{(n)}\right\}} % {y1, ..., yn}, input data
\newcommand{\yvec}{\left(y^{(1)}, \hdots, y^{(n)}\right)^\top} % (y1, ..., yn), vector of outcomes
\newcommand{\greekxi}{\xi} % Greek letter xi
\renewcommand{\xi}[1][i]{\xv^{(#1)}} % x^i, i-th observed value of x
\newcommand{\yi}[1][i]{y^{(#1)}} % y^i, i-th observed value of y
\newcommand{\xivec}{\left(x^{(i)}_1, \ldots, x^{(i)}_p\right)^\top} % (x1^i, ..., xp^i), i-th observation vector
\newcommand{\xj}{\xv_j} % x_j, j-th feature
\newcommand{\xjvec}{\left(x^{(1)}_j, \ldots, x^{(n)}_j\right)^\top} % (x^1_j, ..., x^n_j), j-th feature vector
\newcommand{\phiv}{\mathbf{\phi}} % Basis transformation function phi
\newcommand{\phixi}{\mathbf{\phi}^{(i)}} % Basis transformation of xi: phi^i := phi(xi)

%%%%%% ml - models general
\newcommand{\lamv}{\bm{\lambda}} % lambda vector, hyperconfiguration vector
\newcommand{\Lam}{\bm{\Lambda}}	 % Lambda, space of all hpos
% Inducer / Inducing algorithm
\newcommand{\preimageInducer}{\left(\defAllDatasets\right)\times\Lam} % Set of all datasets times the hyperparameter space
\newcommand{\preimageInducerShort}{\allDatasets\times\Lam} % Set of all datasets times the hyperparameter space
% Inducer / Inducing algorithm
\newcommand{\ind}{\mathcal{I}} % Inducer, inducing algorithm, learning algorithm

% continuous prediction function f
\newcommand{\ftrue}{f_{\text{true}}}  % True underlying function (if a statistical model is assumed)
\newcommand{\ftruex}{\ftrue(\xv)} % True underlying function (if a statistical model is assumed)
\newcommand{\fx}{f(\xv)} % f(x), continuous prediction function
\newcommand{\fdomains}{f: \Xspace \rightarrow \R^g} % f with domain and co-domain
\newcommand{\Hspace}{\mathcal{H}} % hypothesis space where f is from
\newcommand{\fbayes}{f^{\ast}} % Bayes-optimal model
\newcommand{\fxbayes}{f^{\ast}(\xv)} % Bayes-optimal model
\newcommand{\fkx}[1][k]{f_{#1}(\xv)} % f_j(x), discriminant component function
\newcommand{\fh}{\hat{f}} % f hat, estimated prediction function
\newcommand{\fxh}{\fh(\xv)} % fhat(x)
\newcommand{\fxt}{f(\xv ~|~ \thetav)} % f(x | theta)
\newcommand{\fxi}{f\left(\xv^{(i)}\right)} % f(x^(i))
\newcommand{\fxih}{\hat{f}\left(\xv^{(i)}\right)} % f(x^(i))
\newcommand{\fxit}{f\left(\xv^{(i)} ~|~ \thetav\right)} % f(x^(i) | theta)
\newcommand{\fhD}{\fh_{\D}} % fhat_D, estimate of f based on D
\newcommand{\fhDtrain}{\fh_{\Dtrain}} % fhat_Dtrain, estimate of f based on D
\newcommand{\fhDnlam}{\fh_{\Dn, \lamv}} %model learned on Dn with hp lambda
\newcommand{\fhDlam}{\fh_{\D, \lamv}} %model learned on D with hp lambda
\newcommand{\fhDnlams}{\fh_{\Dn, \lamv^\ast}} %model learned on Dn with optimal hp lambda
\newcommand{\fhDlams}{\fh_{\D, \lamv^\ast}} %model learned on D with optimal hp lambda

% discrete prediction function h
\newcommand{\hx}{h(\xv)} % h(x), discrete prediction function
\newcommand{\hh}{\hat{h}} % h hat
\newcommand{\hxh}{\hat{h}(\xv)} % hhat(x)
\newcommand{\hxt}{h(\xv | \thetav)} % h(x | theta)
\newcommand{\hxi}{h\left(\xi\right)} % h(x^(i))
\newcommand{\hxit}{h\left(\xi ~|~ \thetav\right)} % h(x^(i) | theta)
\newcommand{\hbayes}{h^{\ast}} % Bayes-optimal classification model
\newcommand{\hxbayes}{h^{\ast}(\xv)} % Bayes-optimal classification model

% yhat
\newcommand{\yh}{\hat{y}} % yhat for prediction of target
\newcommand{\yih}{\hat{y}^{(i)}} % yhat^(i) for prediction of ith targiet
\newcommand{\resi}{\yi- \yih}

% theta
\newcommand{\thetah}{\hat{\theta}} % theta hat
\newcommand{\thetav}{\bm{\theta}} % theta vector
\newcommand{\thetavh}{\bm{\hat\theta}} % theta vector hat
\newcommand{\thetat}[1][t]{\thetav^{[#1]}} % theta^[t] in optimization
\newcommand{\thetatn}[1][t]{\thetav^{[#1 +1]}} % theta^[t+1] in optimization
\newcommand{\thetahDnlam}{\thetavh_{\Dn, \lamv}} %theta learned on Dn with hp lambda
\newcommand{\thetahDlam}{\thetavh_{\D, \lamv}} %theta learned on D with hp lambda
\newcommand{\mint}{\min_{\thetav \in \Theta}} % min problem theta
\newcommand{\argmint}{\argmin_{\thetav \in \Theta}} % argmin theta

% densities + probabilities
% pdf of x
\newcommand{\pdf}{p} % p
\newcommand{\pdfx}{p(\xv)} % p(x)
\newcommand{\pixt}{\pi(\xv~|~ \thetav)} % pi(x|theta), pdf of x given theta
\newcommand{\pixit}[1][i]{\pi\left(\xi[#1] ~|~ \thetav\right)} % pi(x^i|theta), pdf of x given theta
\newcommand{\pixii}[1][i]{\pi\left(\xi[#1]\right)} % pi(x^i), pdf of i-th x

% pdf of (x, y)
\newcommand{\pdfxy}{p(\xv,y)} % p(x, y)
\newcommand{\pdfxyt}{p(\xv, y ~|~ \thetav)} % p(x, y | theta)
\newcommand{\pdfxyit}{p\left(\xi, \yi ~|~ \thetav\right)} % p(x^(i), y^(i) | theta)

% pdf of x given y
\newcommand{\pdfxyk}[1][k]{p(\xv | y= #1)} % p(x | y = k)
\newcommand{\lpdfxyk}[1][k]{\log p(\xv | y= #1)} % log p(x | y = k)
\newcommand{\pdfxiyk}[1][k]{p\left(\xi | y= #1 \right)} % p(x^i | y = k)

% prior probabilities
\newcommand{\pik}[1][k]{\pi_{#1}} % pi_k, prior
\newcommand{\pih}{\hat{\pi}} % pi hat, estimated prior (binary classification)
\newcommand{\pikh}[1][k]{\hat{\pi}_{#1}} % pi_k hat, estimated prior
\newcommand{\lpik}[1][k]{\log \pi_{#1}} % log pi_k, log of the prior
\newcommand{\pit}{\pi(\thetav)} % Prior probability of parameter theta

% posterior probabilities
\newcommand{\post}{\P(y = 1 ~|~ \xv)} % P(y = 1 | x), post. prob for y=1
\newcommand{\postk}[1][k]{\P(y = #1 ~|~ \xv)} % P(y = k | y), post. prob for y=k
\newcommand{\pidomains}{\pi: \Xspace \rightarrow \unitint} % pi with domain and co-domain
\newcommand{\pibayes}{\pi^{\ast}} % Bayes-optimal classification model
\newcommand{\pixbayes}{\pi^{\ast}(\xv)} % Bayes-optimal classification model
\newcommand{\pix}{\pi(\xv)} % pi(x), P(y = 1 | x)
\newcommand{\piv}{\bm{\pi}} % pi, bold, as vector
\newcommand{\pikx}[1][k]{\pi_{#1}(\xv)} % pi_k(x), P(y = k | x)
\newcommand{\pikxt}[1][k]{\pi_{#1}(\xv ~|~ \thetav)} % pi_k(x | theta), P(y = k | x, theta)
\newcommand{\pixh}{\hat \pi(\xv)} % pi(x) hat, P(y = 1 | x) hat
\newcommand{\pikxh}[1][k]{\hat \pi_{#1}(\xv)} % pi_k(x) hat, P(y = k | x) hat
\newcommand{\pixih}{\hat \pi(\xi)} % pi(x^(i)) with hat
\newcommand{\pikxih}[1][k]{\hat \pi_{#1}(\xi)} % pi_k(x^(i)) with hat
\newcommand{\pdfygxt}{p(y ~|~\xv, \thetav)} % p(y | x, theta)
\newcommand{\pdfyigxit}{p\left(\yi ~|~\xi, \thetav\right)} % p(y^i |x^i, theta)
\newcommand{\lpdfygxt}{\log \pdfygxt } % log p(y | x, theta)
\newcommand{\lpdfyigxit}{\log \pdfyigxit} % log p(y^i |x^i, theta)

% probababilistic
\newcommand{\bayesrulek}[1][k]{\frac{\P(\xv | y= #1) \P(y= #1)}{\P(\xv)}} % Bayes rule
\newcommand{\muv}{\bm{\mu}} % expectation vector of Gaussian
\newcommand{\muk}[1][k]{\bm{\mu_{#1}}} % mean vector of class-k Gaussian (discr analysis)
\newcommand{\mukh}[1][k]{\bm{\hat{\mu}_{#1}}} % estimated mean vector of class-k Gaussian (discr analysis)

% residual and margin
\newcommand{\eps}{\epsilon} % residual, stochastic
\newcommand{\epsv}{\bm{\epsilon}} % residual, stochastic, as vector
\newcommand{\epsi}{\epsilon^{(i)}} % epsilon^i, residual, stochastic
\newcommand{\epsh}{\hat{\epsilon}} % residual, estimated
\newcommand{\epsvh}{\hat{\epsv}} % residual, estimated, vector
\newcommand{\yf}{y \fx} % y f(x), margin
\newcommand{\yfi}{\yi \fxi} % y^i f(x^i), margin
\newcommand{\Sigmah}{\hat \Sigma} % estimated covariance matrix
\newcommand{\Sigmahj}{\hat \Sigma_j} % estimated covariance matrix for the j-th class

% ml - loss, risk, likelihood
\newcommand{\Lyf}{L\left(y, f\right)} % L(y, f), loss function
\newcommand{\Lypi}{L\left(y, \pi\right)} % L(y, pi), loss function
\newcommand{\Lxy}{L\left(y, \fx\right)} % L(y, f(x)), loss function
\newcommand{\Lxyi}{L\left(\yi, \fxi\right)} % loss of observation
\newcommand{\Lxyt}{L\left(y, \fxt\right)} % loss with f parameterized
\newcommand{\Lxyit}{L\left(\yi, \fxit\right)} % loss of observation with f parameterized
\newcommand{\Lxym}{L\left(\yi, f\left(\bm{\tilde{x}}^{(i)} ~|~ \thetav\right)\right)} % loss of observation with f parameterized
\newcommand{\Lpixy}{L\left(y, \pix\right)} % loss in classification
\newcommand{\Lpiy}{L\left(y, \pi\right)} % loss in classification
\newcommand{\Lpiv}{L\left(y, \piv\right)} % loss in classification
\newcommand{\Lpixyi}{L\left(\yi, \pixii\right)} % loss of observation in classification
\newcommand{\Lpixyt}{L\left(y, \pixt\right)} % loss with pi parameterized
\newcommand{\Lpixyit}{L\left(\yi, \pixit\right)} % loss of observation with pi parameterized
\newcommand{\Lhy}{L\left(y, h\right)} % L(y, h), loss function on discrete classes
\newcommand{\Lhxy}{L\left(y, \hx\right)} % L(y, h(x)), loss function on discrete classes
\newcommand{\Lr}{L\left(r\right)} % L(r), loss defined on residual (reg) / margin (classif)
\newcommand{\lone}{|y - \fx|} % L1 loss
\newcommand{\ltwo}{\left(y - \fx\right)^2} % L2 loss
\newcommand{\lbernoullimp}{\ln(1 + \exp(-y \cdot \fx))} % Bernoulli loss for -1, +1 encoding
\newcommand{\lbernoullizo}{- y \cdot \fx + \log(1 + \exp(\fx))} % Bernoulli loss for 0, 1 encoding
\newcommand{\lcrossent}{- y \log \left(\pix\right) - (1 - y) \log \left(1 - \pix\right)} % cross-entropy loss
\newcommand{\lbrier}{\left(\pix - y \right)^2} % Brier score
\newcommand{\risk}{\mathcal{R}} % R, risk
\newcommand{\riskbayes}{\mathcal{R}^\ast}
\newcommand{\riskf}{\risk(f)} % R(f), risk
\newcommand{\riskdef}{\E_{y|\xv}\left(\Lxy \right)} % risk def (expected loss)
\newcommand{\riskt}{\mathcal{R}(\thetav)} % R(theta), risk
\newcommand{\riske}{\mathcal{R}_{\text{emp}}} % R_emp, empirical risk w/o factor 1 / n
\newcommand{\riskeb}{\bar{\mathcal{R}}_{\text{emp}}} % R_emp, empirical risk w/ factor 1 / n
\newcommand{\riskef}{\riske(f)} % R_emp(f)
\newcommand{\risket}{\mathcal{R}_{\text{emp}}(\thetav)} % R_emp(theta)
\newcommand{\riskr}{\mathcal{R}_{\text{reg}}} % R_reg, regularized risk
\newcommand{\riskrt}{\mathcal{R}_{\text{reg}}(\thetav)} % R_reg(theta)
\newcommand{\riskrf}{\riskr(f)} % R_reg(f)
\newcommand{\riskrth}{\hat{\mathcal{R}}_{\text{reg}}(\thetav)} % hat R_reg(theta)
\newcommand{\risketh}{\hat{\mathcal{R}}_{\text{emp}}(\thetav)} % hat R_emp(theta)
\newcommand{\LL}{\mathcal{L}} % L, likelihood
\newcommand{\LLt}{\mathcal{L}(\thetav)} % L(theta), likelihood
\newcommand{\LLtx}{\mathcal{L}(\thetav | \xv)} % L(theta|x), likelihood
\newcommand{\logl}{\ell} % l, log-likelihood
\newcommand{\loglt}{\logl(\thetav)} % l(theta), log-likelihood
\newcommand{\logltx}{\logl(\thetav | \xv)} % l(theta|x), log-likelihood
\newcommand{\errtrain}{\text{err}_{\text{train}}} % training error
\newcommand{\errtest}{\text{err}_{\text{test}}} % test error
\newcommand{\errexp}{\overline{\text{err}_{\text{test}}}} % avg training error

% lm
\newcommand{\thx}{\thetav^\top \xv} % linear model
\newcommand{\olsest}{(\Xmat^\top \Xmat)^{-1} \Xmat^\top \yv} % OLS estimator in LM

%%%%%% perturbed data
\usepackage{xifthen}
\newcommand{\pert}[3]{\ifthenelse{\equal{#2}{}}{\tilde{#1}}{\ifthenelse{\equal{#3}{}}{\tilde{#1}^{#2}}{\tilde{#1}^{#2|#3}}}}	% command to express that for #1 the subset #2 was perturbed given subset #3

%%%%%% marginalized functions
\newcommand{\fj}{f_j} % marginal function f_j, depending on feature j
\newcommand{\fnj}{f_{-j}} % marginal function f_{-j}, depending on all features but j
\newcommand{\fS}{f_S} % marginal function f_S depending on feature set S
\newcommand{\fnj}{f_{-S}} % marginal function f_{-S}, depending on all features but features in S
\newcommand{\fC}{f_C} % marginal function f_C depending on feature set C
\newcommand{\fhj}{\fh_j} % marginal function fh_j, depending on feature j
\newcommand{\fhnj}{\fh_{-j}} % marginal function fh_{-j}, depending on all features but j
\newcommand{\fhS}{\fh_S} % marginal function fh_S depending on feature set S
\newcommand{\fhnS}{\fh_{-S}} % marginal function fh_{-S}, depending on all features but features in S
\newcommand{\fhC}{\fh_C} % marginal function fh_C depending on feature set C
\newcommand{\XSmat}{\Xmat_S} % Design matrix subset
\newcommand{\XnSmat}{\Xmat__{-S}} % Design matrix subset without S
\newcommand{\XCmat}{\Xmat_C} % Design matrix subset
\newcommand{\Xnj}{\Xmat_{-j}} % Design matrix subset -j = {1, .., j-1, j+1, ..., p}

%%%%% Shapley values
\newcommand{\Scupj}{S \cup \{j\}} % coalition S but without player j
\newcommand{\Scupk}{S \cup \{k\}} % coalition S but without player k
\newcommand{\SsubP}{S \subseteq P} % coalition S subset of P
\newcommand{\SsubPnoj}{\SsubP \setminus \{j\}} % coalition S subset of P without player j
\newcommand{\SsubPnojk}{\SsubP \setminus \{j,k\}} % coalition S subset of P without player k
\newcommand{\phiij}{\hat{\phi}_j^{(i)}} % Shapley value for feature j and observation i

%%%%% LIME 
\newcommand{\Gspace}{\mathcal{G}} % Hypothesis space for surrogate model
\newcommand{\neigh}{\phi_{\xv}} % Proximity measure
\newcommand{\zv}{\mathbf{z}} % Sampled datapoints for surrogate
\newcommand{\Zspace}{\mathcal{Z}} % Space of sampled datapoints
\newcommand{\Gower}{d_G} % Gower distance


\newcommand{\indep}{\perp}
\newcommand{\Ind}{\mathcal{I}}

%new environments
\newenvironment{vbframe}  %frame with breaks and verbatim
{
\begin{frame}%[containsverbatim,allowframebreaks]
}
{
\end{frame}
}

% \newenvironment{vframe}  %frame with verbatim without breaks (to avoid numbering one slided frames)
% {
%  \begin{frame}[containsverbatim]
% }
% {
% \end{frame}
% }

\newenvironment{blocki}[1]   % itemize block
{
 \begin{block}{#1}\begin{itemize}
}
{
\end{itemize}\end{block}
}

\newenvironment{fragileframe}[2]{  %fragile frame with framebreaks
\begin{frame}[allowframebreaks, fragile, environment = fragileframe]
\frametitle{#1}
#2}
{\end{frame}}


\newcommand{\myframe}[2]{  %short for frame with framebreaks
\begin{frame}[allowframebreaks]
\frametitle{#1}
#2
\end{frame}}

\newcommand{\remark}[1]{
  \textbf{Remark:} #1
}


\newenvironment{deleteframe}
{
\begingroup
\usebackgroundtemplate{\includegraphics[width=\paperwidth,height=\paperheight]{../style/color/red.png}}
 \begin{frame}
}
{
\end{frame}
\endgroup
}
\newenvironment{simplifyframe}
{
\begingroup
\usebackgroundtemplate{\includegraphics[width=\paperwidth,height=\paperheight]{../style/color/yellow.png}}
 \begin{frame}
}
{
\end{frame}
\endgroup
}\newenvironment{draftframe}
{
\begingroup
\usebackgroundtemplate{\includegraphics[width=\paperwidth,height=\paperheight]{../style/color/green.jpg}}
 \begin{frame}
}
{
\end{frame}
\endgroup
}
% https://tex.stackexchange.com/a/261480: textcolor that works in mathmode
\makeatletter
\renewcommand*{\@textcolor}[3]{%
  \protect\leavevmode
  \begingroup
    \color#1{#2}#3%
  \endgroup
}
\makeatother

% \makeatletter
% %\newcommand\notsotiny{\@setfontsize\notsotiny\@vipt\@viipt}
% \newcommand\notsotiny{\@setfontsize\notsotiny{6.31415}{7.1828}}
% \makeatother

\providecommand{\tightlist}{%
  \setlength{\itemsep}{0pt}\setlength{\parskip}{0pt}}

%\setbeamerfont{footnote}{size=\tiny}
\usepackage[hang,flushmargin]{footmisc}
\renewcommand*{\footnotelayout}{\tiny}
\renewcommand*{\thefootnote}{} %\fnsymbol{footnote}

% https://tex.stackexchange.com/questions/638616/beamer-frametitle-always-uppercase-for-text-but-not-for-href-or-url
\usepackage{textcase}

% https://stackoverflow.com/questions/377301/is-there-a-latex-command-to-make-text-all-lower-caps
%\usepackage[overload]{textcase}
%\setbeamertemplate{frametitle}{\MakeTextUppercase{\insertframetitle}}
%\setbeamertemplate{frametitle}{\expandafter\uppercase\expandafter\insertframetitle}

\setbeamertemplate{frametitle}{\MakeTextUppercase{\insertframetitle}}
\newcommand{\citebutton}[2]{%
\NoCaseChange{\resizebox{!}{9pt}{\protect\beamergotobutton{\href{#2}{#1}}}}%
}

%\setbeamertemplate{frametitle}{\expandafter\uppercase\expandafter\insertframetitle}
%\newcommand{\citebutton}[2]{%
%\href{#2}{\footnotesize\color{black!60}[#1]}
%\resizebox{!}{9pt}{\beamergotobutton{\lowercase{\href{#2}}{#1}}}% \MakeLowercase
%}
%\newcommand{\lit}[2]{\href{#2}{\footnotesize\color{black!60}[#1]}}

\newcommand{\lit}[2]{\href{#2}{\footnotesize\color{black!60}[#1]}}

\let\code=\texttt
\let\proglang=\textsf

\setkeys{Gin}{width=0.9\textwidth}

% https://tex.stackexchange.com/questions/30720/footnote-without-a-marker
% \makeatletter
% \def\blfootnote{\gdef\@thefnmark{}\@footnotetext}
% \makeatother

% https://tex.stackexchange.com/questions/357717/beamer-allowframebreaks-option-and-vertical-spacing-when-using-lists-itemize
% \setbeamertemplate{frametitle continuation}{%
%     (\insertcontinuationcount)%
%     \ifnum\insertcontinuationcount>1%
%     \vspace*{\topsep}%
%     \else%
%     %
%     \fi%
% }


  % math spaces
\ifdefined\N                                                                
\renewcommand{\N}{\mathds{N}} % N, naturals
\else \newcommand{\N}{\mathds{N}} \fi 
\newcommand{\Z}{\mathds{Z}} % Z, integers
\newcommand{\Q}{\mathds{Q}} % Q, rationals
\newcommand{\R}{\mathds{R}} % R, reals
\ifdefined\C 
  \renewcommand{\C}{\mathds{C}} % C, complex
\else \newcommand{\C}{\mathds{C}} \fi
\newcommand{\continuous}{\mathcal{C}} % C, space of continuous functions
\newcommand{\M}{\mathcal{M}} % machine numbers
\newcommand{\epsm}{\epsilon_m} % maximum error

% counting / finite sets
\newcommand{\setzo}{\{0, 1\}} % set 0, 1
\newcommand{\setmp}{\{-1, +1\}} % set -1, 1
\newcommand{\unitint}{[0, 1]} % unit interval

% basic math stuff
\newcommand{\xt}{\tilde x} % x tilde
\newcommand{\argmax}{\operatorname{arg\,max}} % argmax
\newcommand{\argmin}{\operatorname{arg\,min}} % argmin
\newcommand{\argminlim}{\mathop{\mathrm{arg\,min}}\limits} % argmax with limits
\newcommand{\argmaxlim}{\mathop{\mathrm{arg\,max}}\limits} % argmin with limits  
\newcommand{\sign}{\operatorname{sign}} % sign, signum
\newcommand{\I}{\mathbb{I}} % I, indicator
\newcommand{\order}{\mathcal{O}} % O, order
\newcommand{\pd}[2]{\frac{\partial{#1}}{\partial #2}} % partial derivative
\newcommand{\floorlr}[1]{\left\lfloor #1 \right\rfloor} % floor
\newcommand{\ceillr}[1]{\left\lceil #1 \right\rceil} % ceiling

% sums and products
\newcommand{\sumin}{\sum\limits_{i=1}^n} % summation from i=1 to n
\newcommand{\sumim}{\sum\limits_{i=1}^m} % summation from i=1 to m
\newcommand{\sumjn}{\sum\limits_{j=1}^n} % summation from j=1 to p
\newcommand{\sumjp}{\sum\limits_{j=1}^p} % summation from j=1 to p
\newcommand{\sumik}{\sum\limits_{i=1}^k} % summation from i=1 to k
\newcommand{\sumkg}{\sum\limits_{k=1}^g} % summation from k=1 to g
\newcommand{\sumjg}{\sum\limits_{j=1}^g} % summation from j=1 to g
\newcommand{\meanin}{\frac{1}{n} \sum\limits_{i=1}^n} % mean from i=1 to n
\newcommand{\meankg}{\frac{1}{g} \sum\limits_{k=1}^g} % mean from k=1 to g
\newcommand{\prodin}{\prod\limits_{i=1}^n} % product from i=1 to n
\newcommand{\prodkg}{\prod\limits_{k=1}^g} % product from k=1 to g
\newcommand{\prodjp}{\prod\limits_{j=1}^p} % product from j=1 to p

% linear algebra
\newcommand{\one}{\boldsymbol{1}} % 1, unitvector
\newcommand{\zero}{\mathbf{0}} % 0-vector
\newcommand{\id}{\boldsymbol{I}} % I, identity
\newcommand{\diag}{\operatorname{diag}} % diag, diagonal
\newcommand{\trace}{\operatorname{tr}} % tr, trace
\newcommand{\spn}{\operatorname{span}} % span
\newcommand{\scp}[2]{\left\langle #1, #2 \right\rangle} % <.,.>, scalarproduct
\newcommand{\mat}[1]{\begin{pmatrix} #1 \end{pmatrix}} % short pmatrix command
\newcommand{\Amat}{\mathbf{A}} % matrix A
\newcommand{\Deltab}{\mathbf{\Delta}} % error term for vectors

% basic probability + stats
\renewcommand{\P}{\mathds{P}} % P, probability
\newcommand{\E}{\mathds{E}} % E, expectation
\newcommand{\var}{\mathsf{Var}} % Var, variance
\newcommand{\cov}{\mathsf{Cov}} % Cov, covariance
\newcommand{\corr}{\mathsf{Corr}} % Corr, correlation
\newcommand{\normal}{\mathcal{N}} % N of the normal distribution
\newcommand{\iid}{\overset{i.i.d}{\sim}} % dist with i.i.d superscript
\newcommand{\distas}[1]{\overset{#1}{\sim}} % ... is distributed as ...
  % machine learning
\newcommand{\Xspace}{\mathcal{X}} % X, input space
\newcommand{\Yspace}{\mathcal{Y}} % Y, output space
\newcommand{\Zspace}{\mathcal{Z}} % Z, space of sampled datapoints
\newcommand{\nset}{\{1, \ldots, n\}} % set from 1 to n
\newcommand{\pset}{\{1, \ldots, p\}} % set from 1 to p
\newcommand{\gset}{\{1, \ldots, g\}} % set from 1 to g
\newcommand{\Pxy}{\mathbb{P}_{xy}} % P_xy
\newcommand{\Exy}{\mathbb{E}_{xy}} % E_xy: Expectation over random variables xy
\newcommand{\xv}{\mathbf{x}} % vector x (bold)
\newcommand{\Xv}{\mathbf{X}} % vector X (random variable, bold)
\newcommand{\xtil}{\tilde{\mathbf{x}}} % vector x-tilde (bold)
\newcommand{\yv}{\mathbf{y}} % vector y (bold)
\newcommand{\Yv}{\mathbf{Y}} % vector Y (random variable, bold)
\newcommand{\xy}{(\xv, y)} % observation (x, y)
\newcommand{\xvec}{\left(x_1, \ldots, x_p\right)^\top} % (x1, ..., xp)
\newcommand{\Xmat}{\mathbf{X}} % Design matrix
\newcommand{\allDatasets}{\mathds{D}} % The set of all datasets
\newcommand{\allDatasetsn}{\mathds{D}_n}  % The set of all datasets of size n
\newcommand{\D}{\mathcal{D}} % D, data
\newcommand{\Dn}{\D_n} % D_n, data of size n
\newcommand{\Dtrain}{\mathcal{D}_{\text{train}}} % D_train, training set
\newcommand{\Dtest}{\mathcal{D}_{\text{test}}} % D_test, test set
\newcommand{\xyi}[1][i]{\left(\xv^{(#1)}, y^{(#1)}\right)} % (x^i, y^i), i-th observation
\newcommand{\Dset}{\left( \xyi[1], \ldots, \xyi[n]\right)} % {(x1,y1)), ..., (xn,yn)}, data
\newcommand{\defAllDatasetsn}{(\Xspace \times \Yspace)^n} % Def. of the set of all datasets of size n
\newcommand{\defAllDatasets}{\bigcup_{n \in \N}(\Xspace \times \Yspace)^n} % Def. of the set of all datasets
\newcommand{\xdat}{\left\{ \xv^{(1)}, \ldots, \xv^{(n)}\right\}} % {x1, ..., xn}, input data
\newcommand{\ydat}{\left\{ \yv^{(1)}, \ldots, \yv^{(n)}\right\}} % {y1, ..., yn}, input data
\newcommand{\yvec}{\left(y^{(1)}, \hdots, y^{(n)}\right)^\top} % (y1, ..., yn), vector of outcomes
\newcommand{\greekxi}{\xi} % Greek letter xi
\renewcommand{\xi}[1][i]{\xv^{(#1)}} % x^i, i-th observed value of x
\newcommand{\yi}[1][i]{y^{(#1)}} % y^i, i-th observed value of y
\newcommand{\xivec}{\left(x^{(i)}_1, \ldots, x^{(i)}_p\right)^\top} % (x1^i, ..., xp^i), i-th observation vector
\newcommand{\xj}{\xv_j} % x_j, j-th feature
\newcommand{\xjvec}{\left(x^{(1)}_j, \ldots, x^{(n)}_j\right)^\top} % (x^1_j, ..., x^n_j), j-th feature vector
\newcommand{\phiv}{\mathbf{\phi}} % Basis transformation function phi
\newcommand{\phixi}{\mathbf{\phi}^{(i)}} % Basis transformation of xi: phi^i := phi(xi)

%%%%%% ml - models general
\newcommand{\lamv}{\bm{\lambda}} % lambda vector, hyperconfiguration vector
\newcommand{\Lam}{\bm{\Lambda}}	 % Lambda, space of all hpos
% Inducer / Inducing algorithm
\newcommand{\preimageInducer}{\left(\defAllDatasets\right)\times\Lam} % Set of all datasets times the hyperparameter space
\newcommand{\preimageInducerShort}{\allDatasets\times\Lam} % Set of all datasets times the hyperparameter space
% Inducer / Inducing algorithm
\newcommand{\ind}{\mathcal{I}} % Inducer, inducing algorithm, learning algorithm

% continuous prediction function f
\newcommand{\ftrue}{f_{\text{true}}}  % True underlying function (if a statistical model is assumed)
\newcommand{\ftruex}{\ftrue(\xv)} % True underlying function (if a statistical model is assumed)
\newcommand{\fx}{f(\xv)} % f(x), continuous prediction function
\newcommand{\fdomains}{f: \Xspace \rightarrow \R^g} % f with domain and co-domain
\newcommand{\Hspace}{\mathcal{H}} % hypothesis space where f is from
\newcommand{\fbayes}{f^{\ast}} % Bayes-optimal model
\newcommand{\fxbayes}{f^{\ast}(\xv)} % Bayes-optimal model
\newcommand{\fkx}[1][k]{f_{#1}(\xv)} % f_j(x), discriminant component function
\newcommand{\fh}{\hat{f}} % f hat, estimated prediction function
\newcommand{\fxh}{\fh(\xv)} % fhat(x)
\newcommand{\fxt}{f(\xv ~|~ \thetav)} % f(x | theta)
\newcommand{\fxi}{f\left(\xv^{(i)}\right)} % f(x^(i))
\newcommand{\fxih}{\hat{f}\left(\xv^{(i)}\right)} % f(x^(i))
\newcommand{\fxit}{f\left(\xv^{(i)} ~|~ \thetav\right)} % f(x^(i) | theta)
\newcommand{\fhD}{\fh_{\D}} % fhat_D, estimate of f based on D
\newcommand{\fhDtrain}{\fh_{\Dtrain}} % fhat_Dtrain, estimate of f based on D
\newcommand{\fhDnlam}{\fh_{\Dn, \lamv}} %model learned on Dn with hp lambda
\newcommand{\fhDlam}{\fh_{\D, \lamv}} %model learned on D with hp lambda
\newcommand{\fhDnlams}{\fh_{\Dn, \lamv^\ast}} %model learned on Dn with optimal hp lambda
\newcommand{\fhDlams}{\fh_{\D, \lamv^\ast}} %model learned on D with optimal hp lambda

% discrete prediction function h
\newcommand{\hx}{h(\xv)} % h(x), discrete prediction function
\newcommand{\hh}{\hat{h}} % h hat
\newcommand{\hxh}{\hat{h}(\xv)} % hhat(x)
\newcommand{\hxt}{h(\xv | \thetav)} % h(x | theta)
\newcommand{\hxi}{h\left(\xi\right)} % h(x^(i))
\newcommand{\hxit}{h\left(\xi ~|~ \thetav\right)} % h(x^(i) | theta)
\newcommand{\hbayes}{h^{\ast}} % Bayes-optimal classification model
\newcommand{\hxbayes}{h^{\ast}(\xv)} % Bayes-optimal classification model

% yhat
\newcommand{\yh}{\hat{y}} % yhat for prediction of target
\newcommand{\yih}{\hat{y}^{(i)}} % yhat^(i) for prediction of ith targiet
\newcommand{\resi}{\yi- \yih}

% theta
\newcommand{\thetah}{\hat{\theta}} % theta hat
\newcommand{\thetav}{\bm{\theta}} % theta vector
\newcommand{\thetavh}{\bm{\hat\theta}} % theta vector hat
\newcommand{\thetat}[1][t]{\thetav^{[#1]}} % theta^[t] in optimization
\newcommand{\thetatn}[1][t]{\thetav^{[#1 +1]}} % theta^[t+1] in optimization
\newcommand{\thetahDnlam}{\thetavh_{\Dn, \lamv}} %theta learned on Dn with hp lambda
\newcommand{\thetahDlam}{\thetavh_{\D, \lamv}} %theta learned on D with hp lambda
\newcommand{\mint}{\min_{\thetav \in \Theta}} % min problem theta
\newcommand{\argmint}{\argmin_{\thetav \in \Theta}} % argmin theta

% densities + probabilities
% pdf of x
\newcommand{\pdf}{p} % p
\newcommand{\pdfx}{p(\xv)} % p(x)
\newcommand{\pixt}{\pi(\xv~|~ \thetav)} % pi(x|theta), pdf of x given theta
\newcommand{\pixit}[1][i]{\pi\left(\xi[#1] ~|~ \thetav\right)} % pi(x^i|theta), pdf of x given theta
\newcommand{\pixii}[1][i]{\pi\left(\xi[#1]\right)} % pi(x^i), pdf of i-th x

% pdf of (x, y)
\newcommand{\pdfxy}{p(\xv,y)} % p(x, y)
\newcommand{\pdfxyt}{p(\xv, y ~|~ \thetav)} % p(x, y | theta)
\newcommand{\pdfxyit}{p\left(\xi, \yi ~|~ \thetav\right)} % p(x^(i), y^(i) | theta)

% pdf of x given y
\newcommand{\pdfxyk}[1][k]{p(\xv | y= #1)} % p(x | y = k)
\newcommand{\lpdfxyk}[1][k]{\log p(\xv | y= #1)} % log p(x | y = k)
\newcommand{\pdfxiyk}[1][k]{p\left(\xi | y= #1 \right)} % p(x^i | y = k)

% prior probabilities
\newcommand{\pik}[1][k]{\pi_{#1}} % pi_k, prior
\newcommand{\pih}{\hat{\pi}} % pi hat, estimated prior (binary classification)
\newcommand{\pikh}[1][k]{\hat{\pi}_{#1}} % pi_k hat, estimated prior
\newcommand{\lpik}[1][k]{\log \pi_{#1}} % log pi_k, log of the prior
\newcommand{\pit}{\pi(\thetav)} % Prior probability of parameter theta

% posterior probabilities
\newcommand{\post}{\P(y = 1 ~|~ \xv)} % P(y = 1 | x), post. prob for y=1
\newcommand{\postk}[1][k]{\P(y = #1 ~|~ \xv)} % P(y = k | y), post. prob for y=k
\newcommand{\pidomains}{\pi: \Xspace \rightarrow \unitint} % pi with domain and co-domain
\newcommand{\pibayes}{\pi^{\ast}} % Bayes-optimal classification model
\newcommand{\pixbayes}{\pi^{\ast}(\xv)} % Bayes-optimal classification model
\newcommand{\pix}{\pi(\xv)} % pi(x), P(y = 1 | x)
\newcommand{\piv}{\bm{\pi}} % pi, bold, as vector
\newcommand{\pikx}[1][k]{\pi_{#1}(\xv)} % pi_k(x), P(y = k | x)
\newcommand{\pikxt}[1][k]{\pi_{#1}(\xv ~|~ \thetav)} % pi_k(x | theta), P(y = k | x, theta)
\newcommand{\pixh}{\hat \pi(\xv)} % pi(x) hat, P(y = 1 | x) hat
\newcommand{\pikxh}[1][k]{\hat \pi_{#1}(\xv)} % pi_k(x) hat, P(y = k | x) hat
\newcommand{\pixih}{\hat \pi(\xi)} % pi(x^(i)) with hat
\newcommand{\pikxih}[1][k]{\hat \pi_{#1}(\xi)} % pi_k(x^(i)) with hat
\newcommand{\pdfygxt}{p(y ~|~\xv, \thetav)} % p(y | x, theta)
\newcommand{\pdfyigxit}{p\left(\yi ~|~\xi, \thetav\right)} % p(y^i |x^i, theta)
\newcommand{\lpdfygxt}{\log \pdfygxt } % log p(y | x, theta)
\newcommand{\lpdfyigxit}{\log \pdfyigxit} % log p(y^i |x^i, theta)

% probababilistic
\newcommand{\bayesrulek}[1][k]{\frac{\P(\xv | y= #1) \P(y= #1)}{\P(\xv)}} % Bayes rule
\newcommand{\muv}{\bm{\mu}} % expectation vector of Gaussian
\newcommand{\muk}[1][k]{\bm{\mu_{#1}}} % mean vector of class-k Gaussian (discr analysis)
\newcommand{\mukh}[1][k]{\bm{\hat{\mu}_{#1}}} % estimated mean vector of class-k Gaussian (discr analysis)

% residual and margin
\newcommand{\eps}{\epsilon} % residual, stochastic
\newcommand{\epsv}{\bm{\epsilon}} % residual, stochastic, as vector
\newcommand{\epsi}{\epsilon^{(i)}} % epsilon^i, residual, stochastic
\newcommand{\epsh}{\hat{\epsilon}} % residual, estimated
\newcommand{\epsvh}{\hat{\epsv}} % residual, estimated, vector
\newcommand{\yf}{y \fx} % y f(x), margin
\newcommand{\yfi}{\yi \fxi} % y^i f(x^i), margin
\newcommand{\Sigmah}{\hat \Sigma} % estimated covariance matrix
\newcommand{\Sigmahj}{\hat \Sigma_j} % estimated covariance matrix for the j-th class

% ml - loss, risk, likelihood
\newcommand{\Lyf}{L\left(y, f\right)} % L(y, f), loss function
\newcommand{\Lypi}{L\left(y, \pi\right)} % L(y, pi), loss function
\newcommand{\Lxy}{L\left(y, \fx\right)} % L(y, f(x)), loss function
\newcommand{\Lxyi}{L\left(\yi, \fxi\right)} % loss of observation
\newcommand{\Lxyt}{L\left(y, \fxt\right)} % loss with f parameterized
\newcommand{\Lxyit}{L\left(\yi, \fxit\right)} % loss of observation with f parameterized
\newcommand{\Lxym}{L\left(\yi, f\left(\bm{\tilde{x}}^{(i)} ~|~ \thetav\right)\right)} % loss of observation with f parameterized
\newcommand{\Lpixy}{L\left(y, \pix\right)} % loss in classification
\newcommand{\Lpiy}{L\left(y, \pi\right)} % loss in classification
\newcommand{\Lpiv}{L\left(y, \piv\right)} % loss in classification
\newcommand{\Lpixyi}{L\left(\yi, \pixii\right)} % loss of observation in classification
\newcommand{\Lpixyt}{L\left(y, \pixt\right)} % loss with pi parameterized
\newcommand{\Lpixyit}{L\left(\yi, \pixit\right)} % loss of observation with pi parameterized
\newcommand{\Lhy}{L\left(y, h\right)} % L(y, h), loss function on discrete classes
\newcommand{\Lhxy}{L\left(y, \hx\right)} % L(y, h(x)), loss function on discrete classes
\newcommand{\Lr}{L\left(r\right)} % L(r), loss defined on residual (reg) / margin (classif)
\newcommand{\lone}{|y - \fx|} % L1 loss
\newcommand{\ltwo}{\left(y - \fx\right)^2} % L2 loss
\newcommand{\lbernoullimp}{\ln(1 + \exp(-y \cdot \fx))} % Bernoulli loss for -1, +1 encoding
\newcommand{\lbernoullizo}{- y \cdot \fx + \log(1 + \exp(\fx))} % Bernoulli loss for 0, 1 encoding
\newcommand{\lcrossent}{- y \log \left(\pix\right) - (1 - y) \log \left(1 - \pix\right)} % cross-entropy loss
\newcommand{\lbrier}{\left(\pix - y \right)^2} % Brier score
\newcommand{\risk}{\mathcal{R}} % R, risk
\newcommand{\riskbayes}{\mathcal{R}^\ast}
\newcommand{\riskf}{\risk(f)} % R(f), risk
\newcommand{\riskdef}{\E_{y|\xv}\left(\Lxy \right)} % risk def (expected loss)
\newcommand{\riskt}{\mathcal{R}(\thetav)} % R(theta), risk
\newcommand{\riske}{\mathcal{R}_{\text{emp}}} % R_emp, empirical risk w/o factor 1 / n
\newcommand{\riskeb}{\bar{\mathcal{R}}_{\text{emp}}} % R_emp, empirical risk w/ factor 1 / n
\newcommand{\riskef}{\riske(f)} % R_emp(f)
\newcommand{\risket}{\mathcal{R}_{\text{emp}}(\thetav)} % R_emp(theta)
\newcommand{\riskr}{\mathcal{R}_{\text{reg}}} % R_reg, regularized risk
\newcommand{\riskrt}{\mathcal{R}_{\text{reg}}(\thetav)} % R_reg(theta)
\newcommand{\riskrf}{\riskr(f)} % R_reg(f)
\newcommand{\riskrth}{\hat{\mathcal{R}}_{\text{reg}}(\thetav)} % hat R_reg(theta)
\newcommand{\risketh}{\hat{\mathcal{R}}_{\text{emp}}(\thetav)} % hat R_emp(theta)
\newcommand{\LL}{\mathcal{L}} % L, likelihood
\newcommand{\LLt}{\mathcal{L}(\thetav)} % L(theta), likelihood
\newcommand{\LLtx}{\mathcal{L}(\thetav | \xv)} % L(theta|x), likelihood
\newcommand{\logl}{\ell} % l, log-likelihood
\newcommand{\loglt}{\logl(\thetav)} % l(theta), log-likelihood
\newcommand{\logltx}{\logl(\thetav | \xv)} % l(theta|x), log-likelihood
\newcommand{\errtrain}{\text{err}_{\text{train}}} % training error
\newcommand{\errtest}{\text{err}_{\text{test}}} % test error
\newcommand{\errexp}{\overline{\text{err}_{\text{test}}}} % avg training error

% lm
\newcommand{\thx}{\thetav^\top \xv} % linear model
\newcommand{\olsest}{(\Xmat^\top \Xmat)^{-1} \Xmat^\top \yv} % OLS estimator in LM

  %%%%%% perturbed data
\usepackage{xifthen}
\newcommand{\pert}[3]{\ifthenelse{\equal{#2}{}}{\tilde{#1}}{\ifthenelse{\equal{#3}{}}{\tilde{#1}^{#2}}{\tilde{#1}^{#2|#3}}}}	% command to express that for #1 the subset #2 was perturbed given subset #3

%%%%%% marginalized functions
\newcommand{\fj}{f_j} % marginal function f_j, depending on feature j
\newcommand{\fnj}{f_{-j}} % marginal function f_{-j}, depending on all features but j
\newcommand{\fS}{f_S} % marginal function f_S depending on feature set S
\newcommand{\fnj}{f_{-S}} % marginal function f_{-S}, depending on all features but features in S
\newcommand{\fC}{f_C} % marginal function f_C depending on feature set C
\newcommand{\fhj}{\fh_j} % marginal function fh_j, depending on feature j
\newcommand{\fhnj}{\fh_{-j}} % marginal function fh_{-j}, depending on all features but j
\newcommand{\fhS}{\fh_S} % marginal function fh_S depending on feature set S
\newcommand{\fhnS}{\fh_{-S}} % marginal function fh_{-S}, depending on all features but features in S
\newcommand{\fhC}{\fh_C} % marginal function fh_C depending on feature set C
\newcommand{\XSmat}{\Xmat_S} % Design matrix subset
\newcommand{\XnSmat}{\Xmat__{-S}} % Design matrix subset without S
\newcommand{\XCmat}{\Xmat_C} % Design matrix subset
\newcommand{\Xnj}{\Xmat_{-j}} % Design matrix subset -j = {1, .., j-1, j+1, ..., p}

%%%%% Shapley values
\newcommand{\Scupj}{S \cup \{j\}} % coalition S but without player j
\newcommand{\Scupk}{S \cup \{k\}} % coalition S but without player k
\newcommand{\SsubP}{S \subseteq P} % coalition S subset of P
\newcommand{\SsubPnoj}{\SsubP \setminus \{j\}} % coalition S subset of P without player j
\newcommand{\SsubPnojk}{\SsubP \setminus \{j,k\}} % coalition S subset of P without player k
\newcommand{\phiij}{\hat{\phi}_j^{(i)}} % Shapley value for feature j and observation i

%%%%% LIME 
\newcommand{\Gspace}{\mathcal{G}} % Hypothesis space for surrogate model
\newcommand{\neigh}{\phi_{\xv}} % Proximity measure
\newcommand{\zv}{\mathbf{z}} % Sampled datapoints for surrogate
\newcommand{\Zspace}{\mathcal{Z}} % Space of sampled datapoints
\newcommand{\Gower}{d_G} % Gower distance


\else
  % ------------------------------------------------------------------
  %  Standalone fallback (approximates the course look)
  % ------------------------------------------------------------------
  \usepackage[english]{babel}
  \usepackage[utf8]{inputenc}
  \usepackage[T1]{fontenc}
  \usepackage[scaled=.92]{helvet}
  \usepackage{amsmath,amssymb,bm,mathtools,csquotes}
  \IfFileExists{dsfont.sty}{\usepackage{dsfont}}{\providecommand{\mathds}{\mathbb}}
  \usepackage{booktabs,multirow,colortbl}
  \usepackage{tikz}
  \usepackage{environ,keyval,pgf}

  \usefonttheme{structurebold}
  \setbeamercolor{structure}{fg=black,bg=white}
  \setbeamercolor{footline}{fg=white,bg=black}
  \setbeamertemplate{navigation symbols}{}
  \setbeamertemplate{itemize items}[ball]
  \setbeamertemplate{blocks}[rounded]
  \setbeamercolor{block title}{fg=black,bg=black!10}
  \setbeamercolor{block body}{fg=black,bg=black!3}
  \setbeamertemplate{frametitle}{%
    \vspace*{0.35em}{\usebeamerfont{frametitle}\large
    \expandafter\uppercase\expandafter{\insertframetitle}\par}%
    \vspace{-0.35em}{\color{black!60}\rule{\textwidth}{0.4pt}}}
  \setbeamertemplate{footline}{%
    \begin{beamercolorbox}[wd=\paperwidth,ht=2.4ex,dp=1ex,leftskip=1em,rightskip=1em]{footline}%
      \tiny\inserttitle: \insertsubtitle\hfill\insertframenumber
    \end{beamercolorbox}}

  % framev: frame with options fs (font size) and align (top/center/bottom)
  \makeatletter
  \define@key{framev}{fs}{\def\framev@fs{#1}}
  \define@key{framev}{align}{\def\framev@align{#1}}
  \NewEnviron{framev}[2][]{%
    \def\framev@fs{}\def\framev@align{center}%
    \setkeys{framev}{#1}%
    \def\framev@va{c}%
    \def\@tmpa{top}\ifx\framev@align\@tmpa\def\framev@va{t}\fi
    \def\@tmpa{t}\ifx\framev@align\@tmpa\def\framev@va{t}\fi
    \def\@tmpa{bottom}\ifx\framev@align\@tmpa\def\framev@va{b}\fi
    \def\@tmpa{b}\ifx\framev@align\@tmpa\def\framev@va{b}\fi
    \begin{frame}[\framev@va]{#2}%
      \ifx\framev@fs\@empty\else\csname\framev@fs\endcsname\fi
      \BODY
    \end{frame}}
  \makeatother

  % two-column layouts (same interface as the course's splitV.sty)
  \newlength{\splitVsep}\setlength{\splitVsep}{5mm}
  \newcommand{\splitVCalcWidths}[1]{%
    \pgfmathsetmacro{\leftcolwidth}{#1-0.5*\splitVsep/\textwidth}%
    \pgfmathsetmacro{\rightcolwidth}{1-#1-0.5*\splitVsep/\textwidth}}
  \newcommand{\splitVBase}[6][0.5]{%
    \splitVCalcWidths{#1}%
    \begin{minipage}{\textwidth}%
      \begin{minipage}[#2]{\leftcolwidth\textwidth}#4\end{minipage}\hfill
      \begin{minipage}[#5]{\rightcolwidth\textwidth}#6\end{minipage}%
    \end{minipage}}
  \newcommand{\splitVTT}[3][0.5]{\splitVBase[#1]{t}{T}{#2}{t}{#3}}
  \newcommand{\splitVCC}[3][0.5]{\splitVBase[#1]{c}{c}{#2}{c}{#3}}
  \let\splitV\splitVCC

  \newcommand{\spacer}[1][1]{\vspace{#1\baselineskip}}
  \newcommand{\image}[2][1]{\includegraphics[width=#1\textwidth]{#2}}
  \newcommand{\imageC}[2][1]{\begin{center}\includegraphics[width=#1\textwidth]{#2}\end{center}}

  % chunk title slide: {chapter}{chunk title}{figure (optional)}{learning goals}
  \newcommand{\titlemeta}[4]{%
    \begin{frame}[plain,noframenumbering]
      \vspace*{1.2em}
      {\footnotesize\color{black!55}\inserttitle\par}
      \vspace{0.9em}
      {\large\bfseries #1\par}
      \vspace{0.25em}
      {\LARGE\bfseries #2\par}
      \vspace{1.6em}
      \begin{block}{Learning goals}
        \begin{itemize}\setlength{\itemsep}{0.3em} #4 \end{itemize}
      \end{block}
    \end{frame}}
  \def\endlecture{}

  % subset of latex-math (slds-lmu) used in this deck
  \providecommand{\xv}{\mathbf{x}}
  \providecommand{\fh}{\hat{f}}
  \providecommand{\ftrue}{f_{\text{true}}}
  \providecommand{\E}{\mathds{E}}
  \renewcommand{\P}{\mathds{P}}
  \providecommand{\var}{\mathsf{Var}}
  \providecommand{\normal}{\mathcal{N}}
  \providecommand{\D}{\mathcal{D}}
  \providecommand{\Xspace}{\mathcal{X}}
  \providecommand{\Hspace}{\mathcal{H}}
  \providecommand{\R}{\mathds{R}}
  \providecommand{\I}{\mathbb{I}}
  \providecommand{\thetav}{\bm{\theta}}
  \providecommand{\lamv}{\bm{\lambda}}
  \providecommand{\greekxi}{\xi}
  \renewcommand{\xi}[1][i]{\xv^{(#1)}}
  \providecommand{\yi}[1][i]{y^{(#1)}}
  \providecommand{\fhice}[1]{\fh_{#1,ICE}}
  \providecommand{\iid}{\overset{i.i.d}{\sim}}
\fi

% =====================================================================
%  Session-specific definitions (both modes)
% =====================================================================
\usepackage{amssymb}
\usepackage{pgfplots}
\pgfplotsset{compat=1.18}
\usepgfplotslibrary{fillbetween}
\usetikzlibrary{arrows.meta,positioning,calc,fit,decorations.pathreplacing,shapes.geometric}

% list environments: use the course versions if they exist
\makeatletter
\@ifundefined{itemizeM}{\newenvironment{itemizeM}{\begin{itemize}\setlength{\itemsep}{0.45em}}{\end{itemize}}}{}
\@ifundefined{itemizeS}{\newenvironment{itemizeS}{\begin{itemize}\setlength{\itemsep}{0.15em}}{\end{itemize}}}{}
\makeatother

\definecolor{evblue}{HTML}{4D88FF}   % course logo blue
\definecolor{evorange}{HTML}{D55E00} % second accent (colorblind-safe)
\definecolor{evgray}{gray}{0.45}
\hypersetup{colorlinks=true,urlcolor=evblue!70!black,linkcolor=black}

% PD notation used today (consistent with the PD chunk of the course)
\newcommand{\fpdtrue}[1]{f^{\,\text{true}}_{#1,PD}}  % DGP-PD (ground truth)
\newcommand{\fpdbar}[1]{\bar{f}_{#1,PD}}             % learner-PD

% ---------------------------------------------------------------------
% Citations: \defref{KEY}{Short form}{URL}
%   \citeref{KEY}  -> [Short form], linked
%   \readbtn{KEY}  -> small linked button, for frame titles. Keys are
%                     upper case so that upper-cased frame titles are safe.
% ---------------------------------------------------------------------
\makeatletter
\newcommand{\defref}[3]{%
  \@namedef{evref@lab@#1}{#2}%
  \expandafter\edef\csname evref@url@#1\endcsname{\detokenize{#3}}}
\newcommand{\citeref}[1]{%
  \href{\csname evref@url@#1\endcsname}{[\csname evref@lab@#1\endcsname]}}
\newcommand{\readbtn}[1]{%
  \href{\csname evref@url@#1\endcsname}{\tikz[baseline=(rb.base)]
    \node[fill=evblue!22,rounded corners=1.5pt,inner sep=1.6pt,
          font=\tiny\normalfont\color{black}](rb){\csname evref@lab@#1\endcsname};}}
\makeatother

\newcommand{\ok}{{\color{evblue!80!black}$\checkmark$}}
\newcommand{\okc}{{\color{evblue!80!black}($\checkmark$)}}
\newcommand{\riskmark}{{\color{evorange}\textbf{!}}}

% placeholder for a figure from a paper: [height]{what}{caption}{KEY}
\newcommand{\figplaceholder}[4][0.4]{%
  \begin{center}\fbox{\parbox[c][#1\textheight][c]{0.86\textwidth}{\centering
    \textbf{Insert #2}\\[0.35em]{\scriptsize #3}\\[0.35em]{\scriptsize\citeref{#4}}}}\end{center}}

% ---------------------------------------------------------------------
% References used in this session (all links retrieved while preparing)
% ---------------------------------------------------------------------
\defref{DOSHIVELEZ2017}{Doshi-Velez and Kim 2017}{https://arxiv.org/abs/1702.08608v2}
\defref{NAUTA2023}{Nauta et al. 2023}{https://doi.org/10.1145/3583558}
\defref{JACOVI2020}{Jacovi and Goldberg 2020}{https://aclanthology.org/2020.acl-main.386}
\defref{MOLNAR2023}{Molnar et al. 2023}{https://doi.org/10.1007/978-3-031-44064-9_24}
\defref{MOLNAR2020}{Molnar et al. 2020}{https://doi.org/10.1007/978-3-030-43823-4_17}
\defref{MOLNAR2022}{Molnar et al. 2022}{https://doi.org/10.1007/978-3-031-04083-2_4}
\defref{MOLNARBOOK}{Molnar, iML book, Ch.~32}{https://christophm.github.io/interpretable-ml-book/evaluation.html}
\defref{MOOSBAUER2021}{Moosbauer et al. 2021}{https://papers.nips.cc/paper/2021/hash/12ced2db6f0193dda91ba86224ea1cd8-Abstract.html}
\defref{FRIEDMAN2008}{Friedman and Popescu 2008}{https://projecteuclid.org/euclid.aoas/1223908046}
\defref{GOLDSTEIN2015}{Goldstein et al. 2015}{https://arxiv.org/abs/1309.6392}
\defref{APLEY2020}{Apley and Zhu 2020}{https://arxiv.org/abs/1612.08468v2}
\defref{HERBINGER2024}{Herbinger et al. 2024}{https://jmlr.org/papers/v25/23-0699.html}
\defref{FREITAS2014}{Freitas 2014}{https://doi.org/10.1145/2594473.2594475}
\defref{LAGE2019}{Lage et al. 2019}{https://ojs.aaai.org/index.php/HCOMP/article/view/5280}
\defref{POURSABZI2021}{Poursabzi-Sangdeh et al. 2021}{https://arxiv.org/abs/1802.07810v3}
\defref{HASE2020}{Hase and Bansal 2020}{https://aclanthology.org/2020.acl-main.491}
\defref{BANSAL2021}{Bansal et al. 2021}{https://arxiv.org/abs/2006.14779}
\defref{KAUR2020}{Kaur et al. 2020}{https://www.microsoft.com/en-us/research/publication/interpreting-interpretability-understanding-data-scientists-use-of-interpretability-tools-for-machine-learning/}
\defref{ADEBAYO2018}{Adebayo et al. 2018}{https://arxiv.org/abs/1810.03292v2}
\defref{BANIECKI2022}{Baniecki et al. 2022}{https://doi.org/10.1007/978-3-031-26409-2_8}
\defref{SEMENOVA2022}{Semenova et al. 2022}{https://arxiv.org/abs/1908.01755v4}
\defref{FISHER2019}{Fisher et al. 2019}{https://www.jmlr.org/papers/v20/18-760.html}
\defref{ALVAREZMELIS2018}{Alvarez-Melis and Jaakkola 2018}{https://arxiv.org/abs/1806.08049v1}
\defref{NADEAU2003}{Nadeau and Bengio 2003}{https://mlanthology.org/mlj/2003/nadeau2003mlj-inference}
\defref{LIU2021}{Liu et al. 2021}{https://arxiv.org/abs/2106.12543v4}
\defref{HOFFMAN2018}{Hoffman et al. 2018}{https://arxiv.org/abs/1812.04608}

\title{Interpretable Machine Learning}
\subtitle{Evaluation of iML}
\date{}

\begin{document}

% =====================================================================
%  OPENING                                              [approx. 4 min]
% =====================================================================

\begin{framev}[align=top]{Evaluation of iML: Today's Session}
\tableofcontents
\end{framev}


\begin{framev}[align=top]{Recap: What We Know So Far}
\textbf{Setting:} model $\fh: \Xspace \to \R$, data $\D = \{(\xi, \yi)\}_{i=1}^n$, features of interest $S$, all other features $-S$.
\spacer[0.5]
\begin{itemizeM}
\item \textbf{ICE:} $\fhice{S}^{(i)}(\xv_S) = \fh(\xv_S, \xv_{-S}^{(i)})$
\item \textbf{PD:} $f_{S,PD}(\xv_S) = \E_{\xv_{-S}}\big(\fh(\xv_S, \xv_{-S})\big)$, \quad
      $\fh_{S,PD}(\xv_S) = \frac{1}{n}\sum_{i=1}^n \fhice{S}^{(i)}(\xv_S)$
\item \textbf{ALE:} $f_{j,ALE}(x_j) = \displaystyle\int_{z_{0,j}}^{x_j} \E_{\xv_{-j}\mid x_j = z_j}\Big(\frac{\partial \fh(z_j, \xv_{-j})}{\partial z_j}\Big)\, dz_j - c_j$
\item \textbf{Interpretable models:} sparse linear models, GAMs, decision trees, rule sets
\end{itemizeM}
{\scriptsize ICE \citeref{GOLDSTEIN2015}, \ ALE \citeref{APLEY2020}\par}
\pause
\spacer[0.5]
\begin{block}{Today's question}
We can \emph{compute} all of these. How do we know whether an explanation is \emph{good}?
\end{block}
\end{framev}


\begin{framev}[align=top]{Notation for Today}
\small
\renewcommand{\arraystretch}{1.25}
\begin{tabular}{@{}>{\raggedright\arraybackslash}p{0.2\textwidth}>{\raggedright\arraybackslash}p{0.75\textwidth}@{}}
\toprule
Symbol & Meaning \\
\midrule
$\ftrue$ & true function of the data-generating process (DGP), $y = \ftrue(\xv) + \epsilon$ \\
$\fh$ & model fitted by a learner on the random training set $\D$ \\
$f_{S,PD}$ & PD function of the model $\fh$ (\emph{model-PD}) \\
$\fpdtrue{S}$ & PD function of $\ftrue$ (\emph{DGP-PD}), the ground truth \\
$\fpdbar{S}$ & \emph{learner-PD}: $\E_{\D}(f_{S,PD})$, averaged over training sets \\
$\fh_{S,PD}$ & Monte Carlo estimate of $f_{S,PD}$ from $m$ background points \\
$e$, $g_e$ & an explanation, and the function $g_e: \Xspace \to \R$ it asserts \\
$H^2_{jk}$ & interaction statistic for features $j$ and $k$ \\
\bottomrule
\end{tabular}
\end{framev}


% =====================================================================
%  CHUNK 1 -- What makes an explanation good?          [approx. 10 min]
%  Type: motivation + conceptual framing
% =====================================================================
\section{What Makes an Explanation Good?}

\titlemeta{Evaluation of iML}{What Makes an Explanation ``Good''?}{}{
\item Distinguish faithfulness to the model from plausibility to humans
\item Name what an explanation can be evaluated against
\item Describe the three levels of evaluation and their trade-offs
}


\begin{framev}[align=top]{Same Model, Two Stories}
\splitVTT[0.44]{
\spacer[0]
\small
\textbf{Data:} $x_1 \sim U(0,1)$, \ $x_2 = x_1 + \delta$,\\ $\delta \sim \normal(0, 0.02^2)$\par
\spacer[0.4]
\textbf{Model}, fits $y = x_1 + \epsilon$ well:\par
\[\fh(\xv) = x_1 + 10\,(x_1 - x_2)^2\]
On the data, $(x_1 - x_2)^2 \approx 0$, so $\fh(\xv) \approx x_1$.\par
\spacer[0.5]
{\centering
\begin{tikzpicture}[scale=1.75]
  \draw[->] (0,0) -- (1.08,0) node[below left, font=\scriptsize] {$x_1$};
  \draw[->] (0,0) -- (0,1.08) node[below left, font=\scriptsize] {$x_2$};
  \foreach \u in {0.04,0.1,...,0.98} \fill[black!70] ({\u},{\u+0.015*sin(900*\u)}) circle (0.35pt);
  \foreach \y in {0.05,0.2,...,0.95} \node[evorange, font=\tiny] at (0.2,\y) {$\times$};
  \node[evorange, font=\tiny, align=left, anchor=west] at (0.32,1.12) {PDP evaluates $\fh$ at\\$x_1 = 0.2$ for all $x_2$};
\end{tikzpicture}\par}
}{
\spacer[0]
\centering
\begin{tikzpicture}
\begin{axis}[width=\linewidth, height=0.52\textheight,
  axis lines=left, xlabel={$x_1$}, ylabel={centered effect},
  xmin=0, xmax=1, ymin=-1.1, ymax=2.4, xtick={0,0.5,1},
  tick label style={font=\scriptsize}, label style={font=\scriptsize},
  legend style={draw=none, font=\scriptsize, at={(0.5,1.0)}, anchor=north}]
  \addplot[very thick, evorange, domain=0:1, samples=60] {(x-0.5) + 10*((x-0.5)^2 - 1/12)};
  \addlegendentry{PDP}
  \addplot[very thick, evblue, domain=0:1, samples=10] {x-0.5};
  \addlegendentry{ALE}
\end{axis}
\end{tikzpicture}
\spacer[0.3]
\raggedright\small
\pause
Same model, same data, two very different curves.\\[0.3em]
\textbf{Which one is correct?}
}
\end{framev}


\begin{framev}[align=top]{Faithfulness vs.\ Plausibility \readbtn{JACOVI2020}}
\begin{definition}
\textbf{Faithfulness:} how accurately an explanation reflects the true reasoning process of the model.\\[0.3em]
\textbf{Plausibility:} how convincing the explanation is to humans.
\end{definition}
\spacer[0.4]
In the example:
\begin{itemizeM}
\item ALE: plausible (``more $x_1$, more $y$'') and faithful to what $\fh$ does \emph{on the data}
\item PDP: faithful to what $\fh$ does \emph{off the data}, where $x_1 \neq x_2$, and implausible
\item[] \pause \hspace*{-1.2em}$\Rightarrow$ Both are correct answers, but to \emph{different questions}
\item Humans can judge plausibility, not faithfulness: faithfulness evaluation should not rely on human judgment \citeref{JACOVI2020}
\end{itemizeM}
\end{framev}


\begin{framev}[align=top]{What Is an Explanation Evaluated Against?}
\centering
\begin{tikzpicture}[
    every node/.style={font=\small},
    ent/.style={draw, rounded corners, thick, minimum width=2.2cm, minimum height=0.9cm, align=center},
    lab/.style={font=\scriptsize, align=center, inner sep=1.5pt}]
  \node[ent, fill=evblue!12] (e) at (0,0) {Explanation $e$};
  \node[ent] (m) at (-3.4,-2.1) {Model $\fh$};
  \node[ent] (d) at (3.4,-2.1) {DGP $\ftrue$};
  \node[ent] (h) at (0,1.8) {Human};
  \draw[very thick, evblue] (e) -- node[lab, above left=1pt] {faithfulness, fidelity\\(Chunks 2, 4)} (m);
  \draw[very thick, evblue] (e) -- node[lab, above right=1pt] {ground truth\\about the world (Chunk 3)} (d);
  \draw[very thick, evblue] (e) -- node[lab, right=3pt] {plausibility, understanding,\\reliance (Chunk 6)} (h);
  \draw[thick, dashed, evgray] (m) -- node[lab, below=2pt] {model quality, e.g., test error} (d);
\end{tikzpicture}
\spacer[0.6]
\raggedright
\begin{itemizeM}
\item Explaining the model $\neq$ explaining the data \citeref{MOLNAR2023}
\item An explanation can be perfectly faithful to a model that is wrong about the world
\end{itemizeM}
\end{framev}


\begin{framev}[align=top]{Why Not Just Look at It?}
\begin{itemizeM}
\item We need interpretability because the problem formalization is \emph{incomplete}: fairness, safety, or scientific insight are not in the loss \citeref{DOSHIVELEZ2017}
\item Incompleteness $\neq$ uncertainty: uncertainty gives quantified variance, incompleteness gives unquantified bias
\item So the training objective cannot tell us whether an explanation is good either
\end{itemizeM}
\pause
\spacer[0.3]
\begin{block}{Current practice}
A review of more than 300 papers that introduce XAI methods found that 1 in 3 evaluate exclusively with anecdotal evidence, and 1 in 5 evaluate with users \citeref{NAUTA2023}.
\end{block}
\spacer[0.3]
Anecdotal evidence: ``look at this nice PDP, it makes sense''.
\end{framev}


\begin{framev}[align=top]{Three Levels of Evaluation \readbtn{DOSHIVELEZ2017}}
\centering
\begin{tikzpicture}[
    lvl/.style={draw, rounded corners, thick, minimum width=7.6cm, minimum height=1.25cm, align=left, text width=7.3cm, font=\small}]
  \node[lvl, fill=evblue!18] (a) at (0,2.9) {\textbf{Application-grounded:} real humans, real task\\[-0.1em]
    {\scriptsize e.g., loan officers decide with and without PDPs}};
  \node[lvl, fill=evblue!10] (b) at (0,1.45) {\textbf{Human-grounded:} lay people, simplified task\\[-0.1em]
    {\scriptsize e.g., crowd workers predict the output of a rule set}};
  \node[lvl, fill=evblue!4] (c) at (0,0) {\textbf{Functionally-grounded:} no humans, formal proxy\\[-0.1em]
    {\scriptsize e.g., number of rules, error of a PDP against ground truth}};
  \draw[-{Stealth[length=2.5mm]}, very thick, evgray] ($(c.south west)+(-0.35,0)$) -- node[sloped, above, font=\scriptsize] {realism, cost} ($(a.north west)+(-0.35,0)$);
  \draw[-{Stealth[length=2.5mm]}, very thick, evgray] ($(a.north east)+(0.35,0)$) -- node[sloped, above, font=\scriptsize] {scalability, control} ($(c.south east)+(0.35,0)$);
\end{tikzpicture}
\spacer[0.6]
\raggedright
\begin{itemizeM}
\item A proxy is only as good as the human evidence that it tracks interpretability
\item Today: functionally-grounded (Chunks 2--5), human-grounded (Chunk 6), pitfalls (Chunk 7)
\end{itemizeM}
\end{framev}


\begin{framev}[align=top]{Summary: What Makes an Explanation Good?}
\begin{itemizeM}
\item An explanation is always evaluated against something: the model (faithfulness), the data-generating process (ground truth), or the human (plausibility, understanding, reliance).
\item Faithful and plausible are different properties. A convincing curve can misrepresent the model, and a faithful curve can look absurd, e.g., a PDP that reports the model's behavior far from the data.
\item Doshi-Velez and Kim distinguish application-, human-, and functionally-grounded evaluation. Moving down the levels trades realism for scalability and control.
\item Showing a few convincing plots is anecdotal evidence, not an evaluation, yet it remains common in the literature.
\end{itemizeM}
\end{framev}


% =====================================================================
%  CHUNK 2 -- Properties of explanations               [approx. 10 min]
%  Type: taxonomy + first formal definitions
% =====================================================================
\section{Properties of Explanations}

\titlemeta{Evaluation of iML}{Properties of Explanations}{}{
\item List the Co-12 properties and group them into content, presentation, and user
\item Define fidelity formally and compute it for a simple model
\item Judge which properties are at risk for interpretable models and feature effects
}


\begin{framev}[align=top]{Explainability Is Not One Number \readbtn{NAUTA2023}}
\scriptsize
\renewcommand{\arraystretch}{1.12}
\begin{tabular}{@{}lll@{}}
\toprule
 & \textbf{Property} & \textbf{Question it asks} \\
\midrule
\multirow{6}{*}{\textbf{Content}}
 & Correctness & Does it describe the model truthfully? (``nothing but the truth'') \\
 & Completeness & How much of the model does it cover? (``the whole truth'') \\
 & Consistency & Identical inputs, identical explanations? \\
 & Continuity & Similar inputs, similar explanations? \\
 & Contrastivity & Does it answer ``why this and not that?'' \\
 & Covariate complexity & Are the features and concepts it uses understandable? \\
\midrule
\multirow{3}{*}{\textbf{Presentation}}
 & Compactness & How large is it? \\
 & Composition & How is it structured and presented? \\
 & Confidence & Does it report (accurate) uncertainty? \\
\midrule
\multirow{3}{*}{\textbf{User}}
 & Context & Does it fit the user's needs and expertise? \\
 & Coherence & Is it consistent with prior knowledge, i.e., plausible? \\
 & Controllability & Can the user interact with it or correct it? \\
\bottomrule
\end{tabular}
\spacer[0.6]
\small
\begin{itemizeM}
\item No single metric covers all twelve: an evaluation should select several \citeref{NAUTA2023}
\end{itemizeM}
\end{framev}


\begin{framev}[align=top]{Correctness and Completeness as Fidelity}
\begin{definition}[Fidelity of a functional explanation]
Let explanation $e$ assert a function $g_e: \Xspace \to \R$. Its fidelity to $\fh$ is
$$R^2(g_e;\fh) = 1 - \frac{\E_{\xv}\big((\fh(\xv) - g_e(\xv))^2\big)}{\var_{\xv}\big(\fh(\xv)\big)} .$$
\end{definition}
\begin{itemizeS}
\item Interpretable model: the explanation \emph{is} the model, $g_e = \fh$, so $R^2 = 1$ by construction
\item Set of 1D PDPs, centered: $g_e(\xv) = \E(\fh) + \sum_j \big(f_{j,PD}(x_j) - \E(f_{j,PD})\big)$
\end{itemizeS}
\pause
\begin{example}
$\fh(\xv) = x_1 x_2$ with $x_1, x_2 \sim U(-1,1)$ independent:
$f_{1,PD}(x_1) = x_1\,\E(x_2) = 0$, likewise $f_{2,PD} = 0$, so $g_e \equiv 0$ and $R^2 = 0$.\\
Both PDPs are flat, yet the model depends on both features.
\end{example}
\end{framev}


\begin{framev}[align=top]{Continuity and Consistency}
\textbf{Continuity:} for a local explanation $e(\xv)$, e.g., the ICE curve of $\xv$ on a grid, written as a vector:
\begin{definition}[Local Lipschitz estimate \citeref{ALVAREZMELIS2018}]
$$\hat L(\xv) = \max_{\xv' \in B_\varepsilon(\xv)} \frac{\| e(\xv) - e(\xv')\|_2}{\|\xv - \xv'\|_2}, \qquad B_\varepsilon(\xv): \text{ball of radius } \varepsilon \text{ around } \xv$$
\end{definition}
\begin{itemizeS}
\item Smooth $\fh$: ICE curves move smoothly with $\xv_{-S}$, so $\hat L$ is bounded
\item Decision tree: the ICE curve (and the decision path) jumps when $\xv_{-S}$ crosses a split
\end{itemizeS}
\pause
\spacer[0.4]
\textbf{Consistency:} the explanation method should not add its own noise
\begin{itemizeS}
\item PDP: a random subsample of $m$ background points gives a different curve per run
\item ALE: the number of intervals $K$ changes the curve (an estimator bias--variance trade-off)
\item[$\Rightarrow$] Report seeds and method hyperparameters, and show the variation
\end{itemizeS}
\end{framev}


\begin{framev}[align=top]{Compactness and Covariate Complexity}
\textbf{Compactness:} size of the explanation relative to human limits
\begin{itemizeS}
\item Sparse linear model: $\|\hat\thetav\|_0$; \ tree: depth, number of leaves; \ rule set: number of rules and conditions
\item Feature effects: $p$ PDPs, plus $\binom{p}{2}$ 2D PDPs if interactions matter
\end{itemizeS}
\begin{example}
$p = 20$ features: 20 one-dimensional PDPs and $\binom{20}{2} = 190$ two-dimensional PDPs. Few users will read 210 plots.
\end{example}
\spacer[0.4]
\textbf{Covariate complexity:} are the features themselves understandable?
\begin{itemizeS}
\item An effect of ``age'' is easy to read; an effect of ``principal component 3'' or of one dummy of a one-hot encoding is not
\end{itemizeS}
\spacer[0.4]
\pause
Compactness and completeness pull in opposite directions (Chunks 4 and 5).
\end{framev}


\begin{framev}[align=top]{Which Properties Are at Risk?}
\centering\small
\renewcommand{\arraystretch}{1.3}
\begin{tabular}{@{}lcccc@{}}
\toprule
 & Correctness & Completeness & Continuity & Compactness \\
\midrule
Sparse linear model & \ok & \ok & \ok & \ok \ (if sparse) \\
Shallow tree & \ok & \ok & \riskmark\ (jumps) & \ok \\
ICE & \okc & \riskmark\ (local only) & \okc & \riskmark\ ($n$ curves) \\
PDP & \okc & \riskmark\ (interactions) & \ok & \ok \\
ALE & \okc & \riskmark\ (interactions) & \ok & \ok \\
\bottomrule
\end{tabular}
\spacer[0.3]
{\scriptsize \ok\ fine \quad \okc\ faithful to $\fh$, but see caveat below \quad \riskmark\ at risk}
\spacer[0.6]
\raggedright
\begin{itemizeS}
\item Interpretable models: explanation = model, so the risk moves to compactness and continuity
\item Feature effects: faithful to $\fh$ up to estimation error, but a PDP may describe $\fh$ far from the data (Chunk 1), and ALE depends on its binning
\item All global feature effects become incomplete under interactions (Chunk 4)
\end{itemizeS}
\end{framev}


\begin{framev}[align=top]{Summary: Properties of Explanations}
\begin{itemizeM}
\item Explanation quality has many facets. The Co-12 properties group them into content, presentation, and user properties, and no single metric covers all of them.
\item Fidelity measures how well the function an explanation asserts matches the model. It is 1 by construction for an interpretable model, and it can be 0 for a set of PDPs even when every feature matters.
\item Continuity and consistency ask whether explanations change for the right reasons, and not because of seeds, bins, or tiny input changes.
\item Interpretable models move the risk from correctness to compactness; feature effects stay compact but risk being incomplete.
\end{itemizeM}
\end{framev}

% =====================================================================
%  CHUNK 3 -- Ground truth for feature effects         [approx. 13 min]
%  Type: formal core, one derivation, worked numeric example
%  (the HPO research slide can be dropped if time is short: -2 min)
% =====================================================================
\section{Ground Truth for Feature Effects}

\titlemeta{Evaluation of iML}{Ground Truth for Feature Effects}{}{
\item Separate the estimand (DGP-PD) from the model-PD and its estimate
\item Decompose the estimation error into bias, model variance, and Monte Carlo variance
\item Construct and interpret pointwise confidence bands for PDPs
}


\begin{framev}[align=top]{Ground Truth Needs a Known Data-Generating Process}
\begin{itemizeS}
\item Real data: there is no ground truth for ``the effect of $x_1$''
\item Simulation: choose $\ftrue$ and $\P(\xv)$ yourself, then compare
\end{itemizeS}
\spacer[0.6]
\centering
\begin{tikzpicture}[
  box/.style={draw, rounded corners, thick, minimum height=0.85cm, minimum width=1.5cm, align=center, font=\small, inner sep=3pt},
  arr/.style={-{Stealth[length=2mm]}, thick},
  lab/.style={font=\scriptsize, midway, above}]
  \node[box, fill=evblue!12] (dgp) at (0,0) {DGP\\$\ftrue$, $\P(\xv)$};
  \node[box] (d) at (2.75,0) {Data\\$\D$};
  \node[box] (m) at (5.2,0) {Model\\$\fh$};
  \node[box] (est) at (7.75,0) {Estimate\\$\fh_{S,PD}$};
  \node[box, fill=evblue!12] (gt) at (3.4,-2.0) {Ground truth\\$\fpdtrue{S}$};
  \node[box, fill=evorange!15] (err) at (7.75,-2.0) {Error\\$(\fh_{S,PD} - \fpdtrue{S})^2$};
  \draw[arr] (dgp) -- node[lab] {sample} (d);
  \draw[arr] (d) -- node[lab] {learner} (m);
  \draw[arr] (m) -- node[lab] {PD} (est);
  \draw[arr] (dgp) |- node[pos=0.25, left, font=\scriptsize, align=right] {analytic or\\large MC} (gt);
  \draw[arr] (gt) -- (err);
  \draw[arr] (est) -- (err);
\end{tikzpicture}
\spacer[0.6]
\raggedright
\begin{itemizeS}
\item Repeat over many draws of $\D$ to get the distribution of the error
\item Synthetic benchmarks with known ground truth exist, e.g., XAI-Bench for feature attribution methods, which we meet later in the course \citeref{LIU2021}
\end{itemizeS}
\end{framev}


\begin{framev}[align=top]{Estimand and Estimators \readbtn{MOLNAR2023}}
\begin{definition}
For a fixed value $\xv_S$:
\vspace{-0.5em}
\begin{align*}
\text{DGP-PD (estimand):} \quad & \fpdtrue{S}(\xv_S) = \E_{\xv_{-S}}\big(\ftrue(\xv_S,\xv_{-S})\big)\\
\text{model-PD:} \quad & f_{S,PD}(\xv_S) = \E_{\xv_{-S}}\big(\fh(\xv_S,\xv_{-S})\big), \quad \fh \text{ fitted on } \D\\
\text{learner-PD:} \quad & \fpdbar{S}(\xv_S) = \E_{\D}\big(f_{S,PD}(\xv_S)\big)\\
\text{estimate:} \quad & \fh_{S,PD}(\xv_S) = \tfrac{1}{m}\textstyle\sum_{i=1}^{m}\fh(\xv_S,\xv_{-S}^{(i)})
\end{align*}
\end{definition}
\begin{itemizeS}
\item Random: the training set $\D$, and the $m$ background points $\xv_{-S}^{(i)} \iid \P(\xv_{-S})$, drawn independently of $\D$
\end{itemizeS}
\spacer[0.6]
\centering
\begin{tikzpicture}[node distance=1.5cm, every node/.style={font=\small},
  arr/.style={{Stealth[length=1.8mm]}-{Stealth[length=1.8mm]}, thick, evgray}]
  \node (t) {$\fpdtrue{S}$};
  \node[right=of t] (b) {$\fpdbar{S}$};
  \node[right=of b] (m) {$f_{S,PD}$};
  \node[right=of m] (e) {$\fh_{S,PD}$};
  \draw[arr] (t) -- node[above, font=\scriptsize] {bias} (b);
  \draw[arr] (b) -- node[above, font=\scriptsize] {model var.} (m);
  \draw[arr] (m) -- node[above, font=\scriptsize] {MC var.} (e);
\end{tikzpicture}
\end{framev}


\begin{framev}[align=top]{Derivation: Three Sources of Error}
Fix $\xv_S$ and write $\hat P = \fh_{S,PD}(\xv_S)$, \ $P = f_{S,PD}(\xv_S)$, \ $\bar P = \fpdbar{S}(\xv_S)$, \ $P^\ast = \fpdtrue{S}(\xv_S)$.
\begin{align*}
\E\big((\hat P - P^\ast)^2\big) &= \big(\E(\hat P) - P^\ast\big)^2 + \var(\hat P) & &{\scriptstyle\text{bias--variance}}\\
\uncover<2->{\E(\hat P) &= \E_{\D}\big(\E(\hat P \mid \D)\big) = \E_{\D}(P) = \bar P & &{\scriptstyle\text{MC unbiased given }\D}}\\
\uncover<3->{\var(\hat P) &= \var_{\D}\big(\E(\hat P \mid \D)\big) + \E_{\D}\big(\var(\hat P \mid \D)\big) & &{\scriptstyle\text{total variance}}}\\
\uncover<4->{&= \var_{\D}(P) + \tfrac{1}{m}\,\E_{\D}\Big(\var_{\xv_{-S}}\big(\fh(\xv_S,\xv_{-S})\big)\Big) & &{\scriptstyle\text{i.i.d.\ background}}}
\end{align*}
\uncover<5->{%
\begin{block}{Result}
\vspace{-0.9em}
\[\text{MSE} = \underbrace{(\bar P - P^\ast)^2}_{\text{bias}^2} + \underbrace{\var_{\D}(P)}_{\text{model variance}} + \underbrace{\tfrac{1}{m}\,\E_{\D}\big(\var_{\xv_{-S}}(\fh(\xv_S,\xv_{-S}))\big)}_{\text{Monte Carlo variance}}\]
\vspace{-0.7em}
\begin{itemizeS}
\item Bias: the learner cannot represent $\ftrue$, or extrapolates where data are sparse
\item Model variance: refit on a new $\D$, get a different PDP
\item MC variance: the spread of the ICE curves at $\xv_S$, divided by $m$
\end{itemizeS}
\end{block}}
\end{framev}


\begin{framev}[align=top]{Worked Example: Putting Numbers In}
\splitVTT[0.47]{
\spacer[0]
\small
\textbf{DGP:} $\ftrue(\xv) = 2x_1 + x_2$,\\ $x_1, x_2 \sim U(0,1)$ independent\par
\spacer[0.4]
\textbf{Model}, with a spurious interaction:
\[\fh(\xv) = 2x_1 + x_2 + 0.5\,x_1x_2\]
\textbf{At $x_1 = 1$:}
\begin{align*}
\fpdtrue{1}(1) &= 2 + \E(x_2) = 2.5\\
f_{1,PD}(1) &= 2 + 1.5\,\E(x_2) = 2.75
\end{align*}
Error of this model's PD: $0.25$
}{
\spacer[0]
\small
\textbf{Estimate} with $m = 3$ background values $x_2 \in \{0.1, 0.4, 0.9\}$:
\begin{align*}
\fh(1, x_2) &= 2 + 1.5\,x_2\\
&\in \{2.15,\ 2.60,\ 3.35\}\\
\fh_{1,PD}(1) &= 8.10/3 = 2.70
\end{align*}
\textbf{MC standard deviation:}
\begin{align*}
\var_{x_2}(2 + 1.5\,x_2) &= 1.5^2 \cdot \tfrac{1}{12} = 0.1875\\
\sqrt{0.1875/3} &= 0.25 \quad (m = 3)\\
\sqrt{0.1875/300} &= 0.025 \quad (m = 300)
\end{align*}
}
\spacer[0.4]
\pause
\begin{block}{Reading}
With $m = 3$, the Monte Carlo noise ($0.25$) is as large as the model's error ($0.25$). Use enough background points before blaming the model.
\end{block}
\end{framev}


\begin{framev}[align=top]{Pointwise Confidence Bands}
\textbf{Model-PD band}, model fixed, captures only Monte Carlo noise:
\[\fh_{S,PD}(\xv_S) \pm z_{1-\alpha/2}\,\frac{\hat\sigma(\xv_S)}{\sqrt{m}}, \quad \hat\sigma^2(\xv_S) = \frac{1}{m-1}\sum_{i=1}^m\big(\fhice{S}^{(i)}(\xv_S) - \fh_{S,PD}(\xv_S)\big)^2\]
\pause
\textbf{Learner-PD band}, includes model variance: refit on $k$ resampled training sets, get PD estimates with mean $\bar P_k$ and sample variance $s^2$ at $\xv_S$.
\begin{itemizeS}
\item Overlapping training sets make $s^2$ too small. Correct it \citeref{NADEAU2003}, as done for PDPs in \citeref{MOLNAR2023}:
\end{itemizeS}
\[\widehat{\var}(\bar P_k) = \Big(\frac{1}{k} + \frac{n_{\text{test}}}{n_{\text{train}}}\Big)\, s^2, \qquad \bar P_k \pm t_{k-1,\,1-\alpha/2}\,\sqrt{\widehat{\var}(\bar P_k)}\]
\pause
\centering
\begin{tikzpicture}
\begin{axis}[width=0.7\textwidth, height=0.27\textheight, axis lines=left,
  xmin=0, xmax=1, ticks=none, xlabel={$x_S$}, ylabel={PD},
  label style={font=\scriptsize},
  legend style={draw=none, font=\tiny, at={(1.02,1)}, anchor=north west}]
  \addplot[name path=up, draw=none, forget plot, domain=0:1, samples=40] {0.9*sin(deg(5.5*x))+0.05 + 0.08 + 1.2*(x-0.5)^2};
  \addplot[name path=lo, draw=none, forget plot, domain=0:1, samples=40] {0.9*sin(deg(5.5*x))+0.05 - 0.08 - 1.2*(x-0.5)^2};
  \addplot[evblue!20] fill between[of=up and lo];
  \addlegendentry{band}
  \addplot[thick, evblue, domain=0:1, samples=40] {0.9*sin(deg(5.5*x))+0.05};
  \addlegendentry{estimate}
  \addplot[thick, dashed, black, domain=0:1, samples=40] {sin(deg(5.5*x))};
  \addlegendentry{ground truth}
\end{axis}
\end{tikzpicture}\\
{\tiny schematic: bands widen where data are sparse}
\end{framev}


\begin{framev}[align=top]{Research Example: Explaining HPO with PDPs \readbtn{MOOSBAUER2021}}
\begin{itemizeS}
\item Explain an HPO run: PDP of a surrogate model of validation error over hyperparameters $\lamv$
\item Bayesian optimization samples densely where performance is good. This \emph{sampling bias} makes the surrogate unreliable elsewhere.
\item A GP surrogate with posterior mean $\hat m$ and covariance $\hat K$ gives the PD and its uncertainty:
\end{itemizeS}
\[\hat m_S(\lamv_S) = \tfrac{1}{n}\textstyle\sum_{i}\hat m(\lamv_S,\lamv_{-S}^{(i)}), \qquad \hat s^2_S(\lamv_S) = \tfrac{1}{n^2}\,\mathbf{1}^\top \hat K(\lamv_S)\,\mathbf{1}\]
\begin{itemizeS}
\item[] \hspace*{-1.2em}{\small in practice approximated by the average posterior variance $\tfrac{1}{n}\sum_i \hat K_{ii}(\lamv_S)$}
\item Splitting the search space into regions with low uncertainty yields reliable regional PDPs
\end{itemizeS}
\figplaceholder[0.17]{Fig.~3 of Moosbauer et al.\ 2021}{PDPs with confidence bands for surrogates trained on BO data and on uniform data, compared with the true PD}{MOOSBAUER2021}
\end{framev}


\begin{framev}[align=top]{Accurate Model, Wrong Effect}
\splitVTT[0.5]{
\spacer[0]
\small
\textbf{DGP:} $y = x_1 + \epsilon$, \ $x_2 = x_1 + \delta$ with small $\delta$\par
\spacer[0.4]
Two models with almost the same test error:
\begin{itemizeS}
\item Model A: $\fh_A(\xv) = x_1$
\item Model B: $\fh_B(\xv) = x_2$
\end{itemizeS}
Their test MSEs differ by $\var(\delta)$ only.\par
\spacer[0.4]
PD of $x_1$:
\begin{itemizeS}
\item A: $x_1 + \text{const}$, the ground truth
\item B: $\E(x_2) = \text{const}$, flat
\end{itemizeS}
}{
\spacer[0]
\centering
\begin{tikzpicture}
\begin{axis}[width=\linewidth, height=0.5\textheight, axis lines=left,
  xmin=0, xmax=1, ymin=-0.65, ymax=0.75, xtick={0,0.5,1}, ytick={-0.5,0,0.5},
  xlabel={$x_1$}, ylabel={centered PD of $x_1$},
  tick label style={font=\scriptsize}, label style={font=\scriptsize},
  legend style={draw=none, font=\scriptsize, at={(0.02,1)}, anchor=north west}]
  \addplot[line width=3pt, black!20, domain=0:1] {x-0.5};
  \addlegendentry{ground truth}
  \addplot[thick, evblue, domain=0:1] {x-0.5};
  \addlegendentry{model A}
  \addplot[thick, evorange, domain=0:1] {0};
  \addlegendentry{model B}
\end{axis}
\end{tikzpicture}
}
\spacer[0.4]
\pause
\begin{block}{Reading}
B's PDP is perfectly faithful to B and wrong about the world. A ground-truth evaluation always evaluates the explanation \emph{and} the model.
\end{block}
\end{framev}


\begin{framev}[align=top]{Summary: Ground Truth for Feature Effects}
\begin{itemizeM}
\item Ground truth for feature effects only exists in simulation: we choose $\ftrue$, compute the DGP-PD, and compare it with the estimated PD.
\item The squared error of a PD estimate decomposes into the learner's bias, the model variance across training sets, and the Monte Carlo variance. The last term shrinks with the number of background points $m$ and grows with the spread of the ICE curves.
\item Model-PD bands capture only Monte Carlo noise. Learner-PD bands from refits also capture model variance, but need a variance correction because the training sets overlap.
\item An accurate model can still have a PDP far from the ground truth, so a ground-truth evaluation always evaluates the explanation and the model together.
\end{itemizeM}
\end{framev}


% =====================================================================
%  CHUNK 4 -- Is the global summary representative?    [approx. 12 min]
%  Type: diagnostic + proof sketch + "what breaks" (H depends on data)
% =====================================================================
\section{Is the Global Summary Representative?}

\titlemeta{Evaluation of iML}{Is the Global Summary Representative?}{}{
\item Quantify how well a PDP represents the underlying ICE curves
\item Define and compute Friedman's H-statistic
\item Explain why H depends on the data distribution, not only on the model
\item Use regional effects when the global summary fails
}


\begin{framev}[align=top]{Averages Can Hide Everything}
\splitVTT[0.47]{
\spacer[0]
\small
$\fh(\xv) = x_1x_2$, \ $x_1, x_2 \sim U(-1,1)$\par
\spacer[0.5]
ICE of observation $i$:
\[\fhice{1}^{(i)}(x_1) = x_1\,x_2^{(i)},\]
a line with slope $x_2^{(i)}$.\par
\spacer[0.5]
PDP:
\[f_{1,PD}(x_1) = x_1\,\E(x_2) = 0\]
\pause
\begin{itemizeS}
\item PDP says: $x_1$ has no effect
\item ICE says: $x_1$ changes every prediction
\end{itemizeS}
}{
\spacer[0]
\centering
\begin{tikzpicture}
\begin{axis}[width=\linewidth, height=0.52\textheight, axis lines=left,
  xmin=-1, xmax=1, ymin=-1, ymax=1, xtick={-1,0,1}, ytick={-1,0,1},
  xlabel={$x_1$}, ylabel={prediction},
  tick label style={font=\scriptsize}, label style={font=\scriptsize},
  legend style={draw=none, font=\scriptsize, at={(0.5,1.02)}, anchor=south}]
  \pgfplotsinvokeforeach{-0.95,-0.7,-0.45,-0.2,0.15,0.4,0.65,0.9}{
    \addplot[black!35, domain=-1:1, forget plot] {#1*x};}
  \addplot[black!35, domain=-1:1] {0.02*x};
  \addlegendentry{ICE}
  \addplot[very thick, evblue, domain=-1:1] {0};
  \addlegendentry{PDP}
\end{axis}
\end{tikzpicture}
}
\par\spacer[0.8]
\pause
$\Rightarrow$ A completeness failure that is invisible in the PDP alone.
\end{framev}


\begin{framev}[align=top]{Heterogeneity of ICE Curves}
Center each ICE curve on the grid $x_S^{(1)},\dots,x_S^{(G)}$ to remove vertical offsets caused by additive effects of other features (cf.\ c-ICE \citeref{GOLDSTEIN2015}):
\[\tilde f^{(i)}(x_S) = \fhice{S}^{(i)}(x_S) - \frac{1}{G}\sum_{g=1}^G \fhice{S}^{(i)}(x_S^{(g)}), \qquad \tilde f_{PD}(x_S) = \frac{1}{n}\sum_{i=1}^n \tilde f^{(i)}(x_S)\]
\begin{definition}[Interaction-related heterogeneity]
\[\mathrm{Het}_S = \frac{1}{nG}\sum_{i=1}^n\sum_{g=1}^G \Big(\tilde f^{(i)}(x_S^{(g)}) - \tilde f_{PD}(x_S^{(g)})\Big)^2\]
\end{definition}
\begin{itemizeS}
\item No interaction with $S$, i.e., $\fh(\xv) = g_S(\xv_S) + h(\xv_{-S})$: all centered curves coincide, so $\mathrm{Het}_S = 0$
\item Example $x_1x_2$ on $[-1,1]^2$, symmetric grid: $\tilde f^{(i)}(x_1) = x_1x_2^{(i)}$, so\\ $\mathrm{Het}_1 \approx \E(x_1^2)\,\var(x_2) = \tfrac13\cdot\tfrac13 = \tfrac19$
\item Recall Chunk 3: the \emph{uncentered} spread drives the Monte Carlo variance
\end{itemizeS}
\end{framev}


\begin{framev}[align=top]{Friedman's H-Statistic \readbtn{FRIEDMAN2008}}
\begin{definition}[Two-way interaction strength]
All PD functions centered to mean zero over the data:
\[H^2_{jk} = \frac{\sum_{i=1}^n\Big(f_{jk,PD}\big(x_j^{(i)},x_k^{(i)}\big) - f_{j,PD}\big(x_j^{(i)}\big) - f_{k,PD}\big(x_k^{(i)}\big)\Big)^2}{\sum_{i=1}^n f_{jk,PD}^2\big(x_j^{(i)},x_k^{(i)}\big)}\]
\end{definition}
\begin{itemizeS}
\item Share of the variance of the joint effect that the sum of the 1D effects does not explain
\item $H^2_{jk} = 0$: no interaction between $j$ and $k$; \ $H^2_{jk} = 1$: the joint effect is pure interaction
\item Interaction of $j$ with all other features:
\[H^2_j = \frac{\sum_{i}\big(\fh(\xi) - f_{j,PD}(x_j^{(i)}) - f_{-j,PD}(\xv_{-j}^{(i)})\big)^2}{\sum_{i}\fh^2(\xi)}\]
\item Cost: each PD value at a data point averages over $n$ points, i.e., $\mathcal{O}(n^2)$ model calls per pair
\end{itemizeS}
\end{framev}


\begin{framev}[align=top]{Proof Sketch: The Two Extremes}
\textbf{(a) Additive in $j$ and $k$:} \ $\fh(\xv) = g_j(x_j) + g_k(x_k) + h(\xv_{-jk})$
\begin{align*}
f_{jk,PD}(x_j,x_k) &= g_j(x_j) + g_k(x_k) + \E(h)\\
f_{j,PD}(x_j) &= g_j(x_j) + \E(g_k) + \E(h)\\
f_{k,PD}(x_k) &= g_k(x_k) + \E(g_j) + \E(h)
\end{align*}
After centering, $f_{jk,PD} = f_{j,PD} + f_{k,PD}$ at every data point, so the numerator is $0$ and $H^2_{jk} = 0$. No independence assumption is needed.\par
\pause
\spacer[0.5]
\textbf{(b) Pure product:} \ $\fh(\xv) = x_jx_k$ with $\E(x_j) = \E(x_k) = 0$, independent
\[f_{j,PD}(x_j) = x_j\,\E(x_k) = 0, \qquad f_{k,PD}(x_k) = 0, \qquad f_{jk,PD}(x_j,x_k) = x_jx_k\]
Numerator $=$ denominator, so $H^2_{jk} = 1$.\par
\pause
\spacer[0.5]
\begin{block}{The assumption that matters}
$\E(x_k) = 0$. If the features are shifted, part of the product moves into the main effects.
\end{block}
\end{framev}


\begin{framev}[align=top]{Numeric Instance: Same Model, Different H}
$\fh(\xv) = x_1x_2$, independent features. \textbf{Case 1:} $x_1,x_2 \sim U(-1,1)$ gives $H^2_{12} = 1$.
\spacer[0.5]
\textbf{Case 2:} $x_1,x_2 \sim U(0,1)$, so $\E(x_k) = \tfrac12$.\\ Centered PD functions:
\begin{align*}
f_{1,PD} &= \tfrac{x_1}{2} - \tfrac14, \qquad f_{2,PD} = \tfrac{x_2}{2} - \tfrac14, \qquad f_{12,PD} = x_1x_2 - \tfrac14\\
\uncover<2->{f_{12,PD} - f_{1,PD} - f_{2,PD} &= x_1x_2 - \tfrac{x_1}{2} - \tfrac{x_2}{2} + \tfrac14 = \big(x_1 - \tfrac12\big)\big(x_2 - \tfrac12\big)}\\
\uncover<3->{H^2_{12} &\to \frac{\E\big((x_1-\frac12)^2\big)\,\E\big((x_2-\frac12)^2\big)}{\var(x_1x_2)}}\\
\uncover<3->{&= \frac{(1/12)^2}{1/9 - 1/16} = \frac{1/144}{7/144} = \frac17 \approx 0.14}
\end{align*}
\uncover<3->{{\small Simulation with $n = 600$: $\hat H^2_{12} = 0.999$ (case 1) and $0.148$ (case 2).}}
\spacer[0.5]
\uncover<4->{%
\begin{block}{Reading}
$H$ is a property of the model \emph{and} the data distribution. Report the data on which it was computed.
\end{block}}
\end{framev}


\begin{framev}[align=top]{When H Is Large: Regional Effects \readbtn{HERBINGER2024}}
\splitVTT[0.52]{
\spacer[0]
\small
Split the feature space into regions in which ICE curves are homogeneous; show one PDP per region.\par
\spacer[0.5]
\textbf{Example} $x_1x_2$ on $[-1,1]^2$, split at $x_2 = 0$:
\begin{align*}
f_{1,PD}(x_1 \mid x_2 \geq 0) &= 0.5\,x_1\\
f_{1,PD}(x_1 \mid x_2 < 0) &= -0.5\,x_1
\end{align*}
since $\E(x_2 \mid x_2 \geq 0) = 0.5$.\par
\spacer[0.3]
Heterogeneity in each region:\\ $\tfrac13\cdot\tfrac{1}{12} = \tfrac{1}{36}$, i.e., $75\%$ less than $\tfrac19$.
}{
\spacer[0]
\centering
\begin{tikzpicture}
\begin{axis}[width=\linewidth, height=0.5\textheight, axis lines=left,
  xmin=-1, xmax=1, ymin=-1, ymax=1, xtick={-1,0,1}, ytick={-1,0,1},
  xlabel={$x_1$}, ylabel={prediction},
  tick label style={font=\scriptsize}, label style={font=\scriptsize},
  legend style={draw=none, font=\tiny, at={(0.5,1.02)}, anchor=south, legend columns=2}]
  \pgfplotsinvokeforeach{0.15,0.4,0.65,0.9}{\addplot[evblue!30, domain=-1:1, forget plot] {#1*x};}
  \pgfplotsinvokeforeach{-0.95,-0.7,-0.45,-0.2}{\addplot[evorange!30, domain=-1:1, forget plot] {#1*x};}
  \addplot[very thick, evblue, domain=-1:1] {0.5*x};
  \addlegendentry{$x_2 \geq 0$}
  \addplot[very thick, evorange, domain=-1:1] {-0.5*x};
  \addlegendentry{$x_2 < 0$}
\end{axis}
\end{tikzpicture}
}
\spacer[0.6]
\begin{itemizeS}
\item GADGET: recursive partitioning that minimizes interaction-related heterogeneity; works for PD, ALE, and SHAP dependence (later in the course) \citeref{HERBINGER2024}
\item Trade-off: regions restore completeness at the cost of compactness
\end{itemizeS}
\end{framev}


\begin{framev}[align=top]{Summary: Is the Global Summary Representative?}
\begin{itemizeM}
\item A global feature effect summarizes ICE curves. It is only complete when the feature does not interact with other features.
\item The heterogeneity of centered ICE curves around the PDP quantifies what the PDP leaves out.
\item Friedman's H-statistic measures the share of a joint effect that the one-dimensional effects cannot explain. It is 0 for models that are additive in the two features, and 1 for a centered pure product.
\item H depends on the data distribution: the same model $x_1x_2$ gives $H^2 = 1$ on $[-1,1]^2$ and $H^2 = 1/7$ on $[0,1]^2$.
\item When heterogeneity is large, regional effects restore completeness at the cost of compactness.
\end{itemizeM}
\end{framev}

% =====================================================================
%  CHUNK 5 -- Measuring how interpretable a model is   [approx. 12 min]
%  Type: functionally-grounded measures + numeric instance + research link
% =====================================================================
\section{Measuring How Interpretable a Model Is}

\titlemeta{Evaluation of iML}{Measuring How Interpretable a Model Is}{}{
\item List complexity proxies for interpretable models and their limits
\item Compute the interaction strength IAS from a first-order ALE model
\item Read accuracy--complexity trade-offs as a Pareto front
}


\begin{framev}[align=top]{Model-Specific Proxies}
\centering\small
\renewcommand{\arraystretch}{1.3}
\begin{tabular}{@{}ll@{}}
\toprule
Model class & Typical proxy for ``interpretable'' \\
\midrule
Sparse linear model & number of nonzero coefficients $\|\hat\thetav\|_0$ \\
GAM & number of shape functions, their wiggliness \\
Decision tree & depth, number of leaves \\
Rule list, decision set & number of rules, conditions per rule \\
\bottomrule
\end{tabular}
\spacer[0.8]
\raggedright
\begin{itemizeM}
\item Functionally grounded: acceptable \emph{if} humans were shown to understand this model class \citeref{DOSHIVELEZ2017}
\item Problem 1: proxies are not comparable across model classes.\\ Is a depth-4 tree simpler than a GAM with 6 smooth terms?
\item Problem 2: is size the right axis at all? (end of this chunk)
\end{itemizeM}
\end{framev}


\begin{framev}[align=top]{A Model-Agnostic View: Functional Decomposition \readbtn{MOLNAR2020}}
\[\fh(\xv) = f_0 + \sum_{j} f_j(x_j) + \sum_{j<k} f_{jk}(x_j,x_k) + \dots + f_{1,\dots,p}(\xv)\]
Use the ALE main effects: \ $f_{ALE1st}(\xv) = f_0 + \sum_{j=1}^p f_{j,ALE}(x_j)$, \ $f_0 = \frac1n\sum_i \fh(\xi)$
\spacer[0.4]
\begin{itemizeS}
\item \textbf{NF:} number of features used (change a feature, check whether any prediction changes)
\item \textbf{IAS:} how much of $\fh$ the main effects miss
\item \textbf{MEC:} how complex the main effects are
\end{itemizeS}
\begin{definition}[Interaction strength]
\[\mathrm{IAS} = \frac{\sum_{i=1}^n\big(\fh(\xi) - f_{ALE1st}(\xi)\big)^2}{\sum_{i=1}^n\big(\fh(\xi) - f_0\big)^2} = 1 - R^2\big(f_{ALE1st};\fh\big)\]
\end{definition}
\begin{itemizeS}
\item The same idea as fidelity in Chunk 2, with ALE instead of PD
\end{itemizeS}
\end{framev}


\begin{framev}[align=top]{IAS: Numeric Instance}
$\fh(\xv) = x_1 + x_2 + x_1x_2$, \ $x_1, x_2 \sim U(-1,1)$ independent
\begin{align*}
\frac{\partial\fh}{\partial x_1} = 1 + x_2 \;&\Rightarrow\; \E(1 + x_2 \mid x_1) = 1 \;\Rightarrow\; f_{1,ALE}(x_1) = x_1\\
&\phantom{\Rightarrow}\;\text{likewise } f_{2,ALE}(x_2) = x_2\\
\uncover<2->{f_{ALE1st}(\xv) &= f_0 + x_1 + x_2, \qquad f_0 = \E(\fh) = 0}\\
\uncover<3->{\mathrm{IAS} &\to \frac{\E\big((x_1x_2)^2\big)}{\E\big((x_1 + x_2 + x_1x_2)^2\big)} = \frac{1/9}{1/3 + 1/3 + 1/9} = \frac17 \approx 0.14}
\end{align*}
\uncover<3->{%
\begin{itemizeS}
\item Cross terms vanish, e.g., $\E(x_1\cdot x_1x_2) = \E(x_1^2)\,\E(x_2) = 0$
\item Simulation with $n = 20{,}000$: $\widehat{\mathrm{IAS}} = 0.144$
\end{itemizeS}}
\spacer[0.4]
\uncover<4->{%
\begin{block}{Reading}
One seventh of the variance of the predictions is invisible in the two ALE plots.
\end{block}}
\end{framev}


\begin{framev}[align=top]{Main Effect Complexity}
Approximate each ALE curve by a piecewise linear $g_j$ with $K_j$ segments. Add segments until $R^2(g_j; f_{j,ALE}) \geq 1 - \varepsilon$ (e.g., $\varepsilon = 0.05$); set slopes to zero where the fit stays good enough.
\begin{definition}
\vspace{-0.6em}
\begin{align*}
\mathrm{MEC}_j &= K_j + \textstyle\sum_{k=1}^{K_j}\I(\beta_{1,k} \neq 0) - 1, \qquad \beta_{1,k}: \text{slope of segment } k\\
\mathrm{MEC} &= \frac{\sum_j V_j\,\mathrm{MEC}_j}{\sum_j V_j}, \qquad V_j = \frac1n\textstyle\sum_{i} f_{j,ALE}\big(x_j^{(i)}\big)^2
\end{align*}
\end{definition}
\centering
\begin{tikzpicture}[x=0.9cm, y=0.9cm, every node/.style={font=\scriptsize}]
  \begin{scope}
    \draw[->, evgray] (0,0) -- (2.2,0); \draw[->, evgray] (0,0) -- (0,1.5);
    \draw[very thick, evblue] (0.1,0.2) -- (2.0,1.3);
    \node at (1.05,-0.35) {linear: $K = 1$, 1 slope};
    \node at (1.05,-0.75) {$\mathrm{MEC}_j = 1$};
  \end{scope}
  \begin{scope}[xshift=3.3cm]
    \draw[->, evgray] (0,0) -- (2.2,0); \draw[->, evgray] (0,0) -- (0,1.5);
    \draw[very thick, evblue] (0.1,0.3) -- (1.05,0.3); \draw[very thick, evblue] (1.05,1.2) -- (2.0,1.2);
    \node at (1.05,-0.35) {step: $K = 2$, 0 slopes};
    \node at (1.05,-0.75) {$\mathrm{MEC}_j = 1$};
  \end{scope}
  \begin{scope}[xshift=6.6cm]
    \draw[->, evgray] (0,0) -- (2.2,0); \draw[->, evgray] (0,0) -- (0,1.5);
    \draw[thin, black!40, domain=0.1:2.0, samples=40] plot (\x, {0.25+1.1*sin(deg(1.55*\x))});
    \draw[very thick, evblue] (0.1,0.42) -- (0.7,1.17) -- (1.4,1.25) -- (2.0,0.45);
    \node at (1.05,-0.35) {curve: $K = 3$, 3 slopes};
    \node at (1.05,-0.75) {$\mathrm{MEC}_j = 5$};
  \end{scope}
\end{tikzpicture}
\spacer[0.2]
\raggedright
\begin{itemizeS}
\item One line or one step costs 1; every further segment or slope costs 1 more
\item Weighting by $V_j$: features with large effects count more
\end{itemizeS}
\end{framev}


\begin{framev}[align=top]{Accuracy vs.\ Interpretability as a Pareto Front}
\splitVTT[0.55]{
\spacer[0]
\centering
\begin{tikzpicture}
\begin{axis}[width=\linewidth, height=0.6\textheight, axis lines=left,
  xmin=0, xmax=10, ymin=0, ymax=10, ticks=none,
  xlabel={complexity, e.g., MEC or IAS}, ylabel={test error},
  label style={font=\scriptsize}]
  \addplot[thick, evblue, mark=*, mark options={fill=evblue}] coordinates {(0.8,8.6) (2.2,5.6) (4.0,3.6) (7.0,2.4) (9.2,2.1)};
  \addplot[only marks, mark=o, thick, evorange] coordinates {(5.5,6.2) (8.0,4.4) (3.2,7.4)};
  \node[font=\tiny, anchor=west] at (axis cs:0.9,8.9) {linear};
  \node[font=\tiny, anchor=west] at (axis cs:2.3,5.95) {GAM};
  \node[font=\tiny, anchor=west] at (axis cs:4.1,3.95) {boosting, shallow};
  \node[font=\tiny, anchor=south] at (axis cs:8.6,2.5) {boosting, deep};
  \node[font=\tiny, anchor=west] at (axis cs:5.6,6.5) {deep tree};
  \node[font=\tiny, evblue, anchor=west] at (axis cs:6.2,0.9) {Pareto front};
\end{axis}
\end{tikzpicture}\\
{\tiny schematic, not real results}
}{
\spacer[0]
\small
\begin{itemizeS}
\item No single best model: each point on the front is a trade-off
\item Hollow points are dominated: another model is both simpler and more accurate
\item Molnar et al.\ demonstrate multi-objective optimization that trades accuracy off against NF, IAS, and MEC \citeref{MOLNAR2020}
\item AutoML view: multi-objective HPO returns the front; the user picks the trade-off
\end{itemizeS}
}
\end{framev}


\begin{framev}[align=top]{Do Proxies Match Humans?}
\begin{itemizeM}
\item Model size as the only criterion of comprehensibility has clear drawbacks; e.g., monotonicity constraints can make models more acceptable to users \citeref{FREITAS2014}
\item Lage et al.\ ran controlled user studies with \emph{decision sets} in two fictional domains, with simulation, verification, and counterfactual tasks \citeref{LAGE2019}:
  \begin{itemizeS}
  \item introducing new concepts (``cognitive chunks'') had the greatest effect on response time
  \item longer explanations slowed people down less; repeated terms had little effect
  \item the effects were consistent across tasks and domains
  \end{itemizeS}
\end{itemizeM}
\pause
\spacer[0.4]
\begin{block}{Reading}
Not all ``size'' is equal. A proxy needs human evidence, which brings us to user studies.
\end{block}
\end{framev}


\begin{framev}[align=top]{Summary: Measuring How Interpretable a Model Is}
\begin{itemizeM}
\item For interpretable models, functionally-grounded evaluation uses model-specific size proxies such as depth, number of rules, or nonzero coefficients. They cannot be compared across model classes.
\item Functional decomposition gives model-agnostic measures: NF counts the features used, IAS measures what the first-order ALE model misses, and MEC counts the line segments needed to describe the main effects.
\item IAS is a fidelity measure: for $x_1 + x_2 + x_1x_2$ on $[-1,1]^2$, the ALE plots miss one seventh of the variance.
\item Accuracy and complexity form a Pareto front, which turns model selection into a multi-objective problem.
\item Size proxies need human validation. In user studies, new concepts cost more time than longer explanations.
\end{itemizeM}
\end{framev}


% =====================================================================
%  CHUNK 6 -- Human-grounded evaluation                [approx. 13 min]
%  Type: study design, case study, sample-size calculation
% =====================================================================
\section{Human-Grounded Evaluation}

\titlemeta{Evaluation of iML}{Human-Grounded Evaluation}{}{
\item Choose a task and an outcome measure for a user study
\item Interpret the results of a classic study on model transparency
\item Estimate the number of participants needed for a two-group comparison
}


\begin{framev}[align=top]{Tasks for User Studies \readbtn{DOSHIVELEZ2017}}
\splitVTT[0.44]{
\spacer[0]
\small
\textbf{Running example:} a decision set for loans (fictional)
\spacer[0.3]
\begin{block}{Model}
\footnotesize
IF income $< 30$k AND debt $> 10$k\\ \quad THEN reject\\
ELSE IF missed payments $\geq 2$\\ \quad THEN reject\\
ELSE accept
\end{block}
\spacer[0.3]
\textbf{Applicant:} income 25k, debt 5k, missed payments 3
}{
\spacer[0]
\small
\begin{itemizeS}
\item \textbf{Forward simulation:} ``What does the model predict?''\\ {\color{evgray}reject, by the second rule}
\item \textbf{Counterfactual simulation:} ``What must change so that the model accepts?''\\ {\color{evgray}at most 1 missed payment}
\item \textbf{Verification:} ``The model says accept. Is that consistent with the rules?''\\ {\color{evgray}no}
\item \textbf{Forced choice:} ``Which of two explanations do you prefer?''\\ {\color{evgray}measures preference, i.e., plausibility}
\end{itemizeS}
}
\spacer[0.6]
\small
Lage et al.\ used simulation, verification, and counterfactual tasks with decision sets \citeref{LAGE2019}.
\end{framev}


\begin{framev}[align=top]{Outcome Measures, Formally}
Participants $u = 1,\dots,U$, items $t = 1,\dots,T$, response $r_{ut}$, model output $\fh(\xv_t)$, true label $y_t$, final decision $d_{ut}$.
\begin{definition}
\vspace{-0.4em}
\begin{align*}
\text{simulation accuracy:} \quad & \mathrm{SimAcc} = \tfrac{1}{UT}\textstyle\sum_{u,t}\I\big(r_{ut} = \fh(\xv_t)\big)\\
\text{reliance, model right:} \quad & A^+ = \P\big(d = \fh(\xv) \mid \fh(\xv) = y\big)\\
\text{reliance, model wrong:} \quad & A^- = \P\big(d = \fh(\xv) \mid \fh(\xv) \neq y\big)
\end{align*}
\end{definition}
\begin{itemizeS}
\item Appropriate reliance: high $A^+$ \emph{and} low $A^-$. Raising both only increases blind trust.
\item Also: response time (often on correct answers only), subjective ratings
\end{itemizeS}
\pause
\begin{itemizeS}
\item Subjective ratings of explanations did not predict how much they actually helped users simulate the model \citeref{HASE2020}
\item Explanations increased the chance that people accepted AI advice, regardless of whether it was correct, i.e., they also raised $A^-$ \citeref{BANSAL2021}
\end{itemizeS}
\end{framev}


\begin{framev}[align=top]{Case Study: Does Transparency Help? \readbtn{POURSABZI2021}}
\small
Predict apartment prices in New York City with the help of a linear model. Pre-registered experiments, $N = 3{,}800$ crowd workers. Identical predictions in all conditions.\par
\spacer[0.9]
\centering
\begin{tikzpicture}[cell/.style={draw, thick, minimum width=3.0cm, minimum height=0.95cm, align=center, font=\scriptsize}]
  \node[cell, fill=evblue!15] at (0,0) {2 features,\\coefficients shown};
  \node[cell] at (3.2,0) {2 features,\\black box};
  \node[cell] at (0,-1.05) {8 features,\\coefficients shown};
  \node[cell] at (3.2,-1.05) {8 features,\\black box};
  \node[font=\scriptsize\bfseries] at (0,0.75) {clear};
  \node[font=\scriptsize\bfseries] at (3.2,0.75) {black box};
  \node[font=\scriptsize, align=right, anchor=east] at (-1.65,0) {bathrooms,\\square footage};
  \node[font=\scriptsize, align=right, anchor=east] at (-1.65,-1.05) {+ 6 more};
\end{tikzpicture}
\spacer[0.5]
\raggedright
\begin{itemizeS}
\item Clear model with 2 features: participants \emph{simulated} the model better
\pause
\item But they did not follow the model more closely when that would have helped
\item Seeing the clear model made it \emph{harder} to detect and correct the model's sizable mistakes; the authors attribute this to information overload
\end{itemizeS}
\pause
\begin{block}{Reading}
Intuition about what is interpretable is a poor guide: test it.
\end{block}
\end{framev}


\begin{framev}[align=top]{Designing the Study}
\begin{itemizeM}
\item \textbf{Between-subjects} (one condition per person) vs.\ \textbf{within-subjects} (all conditions per person: more power, but learning and order effects)
\item \textbf{Baselines:} no explanation, and a random or shuffled explanation. Does the content matter, or only its presence?
\item \textbf{Fictional domains} limit prior knowledge, e.g., recipes for aliens \citeref{LAGE2019}
\item \textbf{Pre-register} hypotheses, outcomes, and exclusion criteria, and report exclusions: in Lage et al., fewer than half of the 900 recruited workers remained after exclusions
\item \textbf{Measure the right construct:} preference measures plausibility, simulation measures understanding, reliance measures decision quality
\end{itemizeM}
\end{framev}


\begin{framev}[align=top]{How Many Participants?}
One binary outcome per participant (e.g., detects the model's error or not), two groups, two-sided test at level $\alpha$, power $1-\beta$, normal approximation:
\begin{definition}
\[n_{\text{per group}} = \frac{\big(z_{1-\alpha/2} + z_{1-\beta}\big)^2\,\big(p_1(1-p_1) + p_2(1-p_2)\big)}{(p_1 - p_2)^2}\]
\end{definition}
\begin{example}
$p_1 = 0.60$ (no explanation), $p_2 = 0.70$ (with explanation), $\alpha = 0.05$, power $0.8$:
\begin{align*}
z_{0.975} = 1.960, \quad z_{0.8} &= 0.842, \quad (1.960 + 0.842)^2 = 7.85\\
n = 7.85 \cdot (0.24 + 0.21)/0.01 &= 353.2 \;\Rightarrow\; 354 \text{ per group, } 708 \text{ in total}
\end{align*}
\end{example}
\pause
\begin{itemizeS}
\item Half the effect, $p_2 = 0.65$: about $1{,}468$ per group, since $n \propto 1/(p_1 - p_2)^2$
\item Several items per participant: responses are correlated, so use mixed-effects models instead of this formula
\end{itemizeS}
\end{framev}


\begin{framev}[align=top]{Summary: Human-Grounded Evaluation}
\begin{itemizeM}
\item Human-grounded evaluation asks people to do something with an explanation: simulate the model, find a counterfactual, verify an output, or choose between explanations.
\item The task determines what is measured: preference measures plausibility, simulation measures understanding, and reliance on correct vs.\ incorrect advice measures decision quality.
\item Intuition is a poor guide. A small, transparent linear model improved simulation but made it harder to catch the model's large errors.
\item Detecting a 10-point improvement in a binary outcome needs about 350 participants per group, so plan the sample size before running a study.
\end{itemizeM}
\end{framev}


% =====================================================================
%  CHUNK 7 -- When evaluation goes wrong               [approx. 12 min]
%  Type: what breaks + connection to current research
% =====================================================================
\section{When Evaluation Goes Wrong}

\titlemeta{Evaluation of iML}{When Evaluation Goes Wrong}{}{
\item Design a randomization sanity check for a feature effect
\item Explain how a PDP can be manipulated without changing the model
\item Define the Rashomon set and its consequence for explanations
}


\begin{framev}[align=top]{Sanity Checks by Randomization \readbtn{ADEBAYO2018}}
\begin{block}{Idea}
An explanation that does not change when the model or the data are destroyed cannot be explaining them.
\end{block}
\spacer[0.4]
\begin{itemizeM}
\item Proposed for saliency maps of image models, which we meet later in the course; the idea transfers to any explanation method
\item \textbf{Model randomization test:} compare the explanation of the trained model with that of an untrained model with random weights
\item \textbf{Data randomization test:} retrain on randomly permuted labels $y^{(\pi(i))}$, compare the explanations
\item Finding: some existing saliency methods are independent of both the model and the data-generating process \citeref{ADEBAYO2018}
\end{itemizeM}
\end{framev}


\begin{framev}[align=top]{A Randomization Test for PDPs}
Let $\fh_\pi$ be trained on permuted labels. On a grid of $G$ points, with centered PD curves $\tilde f_{S,PD}$:
\begin{definition}[our adaptation of the data randomization test]
\[d(\fh, \fh_\pi) = \Big(\frac1G\sum_{g=1}^G\big(\tilde f_{S,PD}(x_S^{(g)}) - \tilde f^{\,\pi}_{S,PD}(x_S^{(g)})\big)^2\Big)^{1/2}\]
Repeat for permutations $\pi_1,\dots,\pi_B$: the curves $\tilde f^{\,\pi_b}_{S,PD}$ form a \emph{null band}.
\end{definition}
\splitVTT[0.55]{
\spacer[0]
\small
\begin{itemizeS}
\item Expected: a flat PDP, up to what the learner fits to noise
\item Effect inside the null band: indistinguishable from fitted noise
\item Same shape as for $\fh$: the shape comes from the learner's inductive bias or from the features, not from $y$
\end{itemizeS}
}{
\spacer[0]
\centering
\begin{tikzpicture}
\begin{axis}[width=\linewidth, height=0.4\textheight, axis lines=left,
  xmin=0, xmax=1, ymin=-0.8, ymax=0.8, ticks=none, xlabel={$x_S$},
  label style={font=\scriptsize},
  legend style={draw=none, font=\tiny, at={(0.02,1)}, anchor=north west}]
  \addplot[name path=nu, draw=none, forget plot, domain=0:1, samples=20] {0.14+0.05*sin(deg(9*x))};
  \addplot[name path=nl, draw=none, forget plot, domain=0:1, samples=20] {-0.14-0.05*sin(deg(7*x))};
  \addplot[black!15] fill between[of=nu and nl];
  \addlegendentry{null band}
  \addplot[very thick, evblue, domain=0:1, samples=40] {0.6*(x-0.5)+0.3*sin(deg(8*x))*x};
  \addlegendentry{PDP of $\fh$}
\end{axis}
\end{tikzpicture}\\
{\tiny schematic}
}
\end{framev}


\begin{framev}[align=top]{Fooling Partial Dependence \readbtn{BANIECKI2022}}
The estimate depends on the background data $X = \{\xv_{-S}^{(i)}\}_{i=1}^m$:
\[\fh_{S,PD}(\xv_S; X) = \tfrac1m\textstyle\sum_{i=1}^m\fh(\xv_S,\xv_{-S}^{(i)})\]
\begin{itemizeS}
\item Attack: keep $\fh$ fixed, poison $X$ so that the curve moves toward a target $t$, e.g.,
\[\min_{X'}\ \textstyle\sum_{g}\big(\fh_{S,PD}(x_S^{(g)}; X') - t(x_S^{(g)})\big)^2 \quad \text{with } X' \text{ close to } X,\]
solved with genetic and gradient-based algorithms
\end{itemizeS}
\begin{example}
$\fh(\xv) = x_1x_2$, so $\fh_{1,PD}(x_1; X) = x_1\cdot\bar x_2(X)$:\\[0.2em]
\begin{tabular}{@{}ll@{}}
background $x_2 \in \{0.2, 0.5, 0.8\}$: & slope $+0.5$\\
background $x_2 \in \{-0.2, -0.5, -0.8\}$: & slope $-0.5$
\end{tabular}\\[0.2em]
Same model, opposite story. For audits: fix, document, and check the background data.
\end{example}
\end{framev}


\begin{framev}[align=top]{The Rashomon Set \readbtn{SEMENOVA2022}}
\begin{definition}[Rashomon set]
\[\hat R(\Hspace,\theta) = \big\{f\in\Hspace : \hat L(f) \leq \hat L(\fh) + \theta\big\},\]
where $\hat L$ is the empirical loss, $\fh$ minimizes $\hat L$ over $\Hspace$, and $\theta \geq 0$.
\end{definition}
\spacer[0.4]
\begin{itemizeM}
\item Chunk 3 again: $x_2 \approx x_1$. Then $f_A = x_1$, $f_B = x_2$, and $f_C = (x_1 + x_2)/2$ all lie in $\hat R$ for a small $\theta$
\item Their PDPs of $x_1$ have slopes $1$, $0$, and $0.5$
\item The data do not identify ``the'' effect of $x_1$; each single model tells one story
\item Fisher et al.\ report the range of variable importance over all well-performing models of a class \citeref{FISHER2019}; the same idea works for feature effects
\end{itemizeM}
\end{framev}


\begin{framev}[align=top]{Evaluating Across the Rashomon Set}
\begin{definition}[PD range, our formulation]
\[\Big[\min_{f\in\hat R} f_{S,PD}(\xv_S),\ \max_{f\in\hat R} f_{S,PD}(\xv_S)\Big]\]
\end{definition}
\splitVTT[0.5]{
\spacer[0]
\small
\begin{itemizeS}
\item In practice, approximate $\hat R$ by many good models, e.g., all configurations of an HPO run within $\theta$ of the best one
\item Narrow range: the effect is a property of the data
\item Wide range: it is a property of the chosen model
\end{itemizeS}
}{
\spacer[0]
\centering
\begin{tikzpicture}
\begin{axis}[width=\linewidth, height=0.45\textheight, axis lines=left,
  xmin=0, xmax=1, ymin=-0.7, ymax=0.8, ticks=none, xlabel={$x_S$}, ylabel={PD},
  label style={font=\scriptsize}]
  \pgfplotsinvokeforeach{0.2,0.4,0.6,0.8,1.0,1.2}{
    \addplot[black!30, domain=0:1, samples=30] {#1*(x-0.5) + 0.12*sin(deg(6*x))*(#1-0.6)};}
  \addplot[very thick, evblue, domain=0:1, samples=30] {0.6*(x-0.5)};
  \node[font=\tiny, evblue, anchor=west] at (axis cs:0.55,-0.45) {chosen model};
\end{axis}
\end{tikzpicture}\\
{\tiny schematic: PDPs of six near-optimal models}
}
\end{framev}


\begin{framev}[align=top]{Pitfalls and the Checks from Today \readbtn{MOLNAR2022}}
\centering\small
\renewcommand{\arraystretch}{1.25}
\begin{tabular}{@{}>{\raggedright\arraybackslash}p{0.43\textwidth}>{\raggedright\arraybackslash}p{0.52\textwidth}@{}}
\toprule
Pitfall & Check from today \\
\midrule
Interpreting a model that does not generalize & report test error before explaining \\
Ignoring feature dependence & compare PDP with ALE; simulate a ground truth (Ch.~1, 3) \\
Ignoring interactions & ICE heterogeneity, H-statistic, IAS (Ch.~4, 5) \\
Ignoring uncertainty & confidence bands, refits (Ch.~3) \\
Unjustified causal interpretation & separate model-PD from DGP-PD (Ch.~3) \\
\bottomrule
\end{tabular}
\spacer[0.7]
\raggedright
\begin{itemizeS}
\item People misuse tools, too: in a contextual inquiry (11) and a survey (197), data scientists over-trusted and misused GAM and SHAP visualizations, and few could explain them accurately \citeref{KAUR2020}
\item So the evaluation of a tool includes how people use it (Ch.~6)
\end{itemizeS}
\end{framev}


\begin{framev}[align=top]{Summary: When Evaluation Goes Wrong}
\begin{itemizeM}
\item Randomization tests check that an explanation depends on the model and on the data. For PDPs, retraining on permuted labels yields a null band against which real effects can be compared.
\item A PDP depends on its background data, so it can be manipulated without touching the model. Audits must control that data.
\item Many models fit the data almost equally well (the Rashomon set), and their feature effects can disagree. The range across the set shows whether an effect is about the data or about one particular model.
\item People misuse explanation tools, too, so evaluating a tool includes evaluating how it is used.
\end{itemizeM}
\end{framev}


% =====================================================================
%  CLOSING                                              [approx. 3 min]
% =====================================================================

\begin{framev}[align=top]{Take-Home: Evaluating an Explanation in Practice}
\begin{enumerate}\setlength{\itemsep}{0.5em}
\item \textbf{Say what you evaluate against:} the model, the data-generating process, or the human.
\item \textbf{Check the model first:} test error, and where the data actually are.
\item \textbf{Quantify uncertainty:} confidence bands, refits, seeds, and estimator settings such as $m$ or the number of ALE intervals.
\item \textbf{Check completeness:} ICE heterogeneity, H, or IAS before trusting one-dimensional summaries.
\item \textbf{Run sanity checks:} permuted labels, and alternative models from the Rashomon set.
\item \textbf{If humans use it, test with humans:} task, outcome measure, baseline, sample size.
\end{enumerate}
\end{framev}


\begin{framev}[align=top]{Further Reading}
\begin{itemizeM}
\item Molnar, \emph{Interpretable Machine Learning}, chapter on evaluation of interpretability methods \citeref{MOLNARBOOK}
\item Nauta et al.: systematic review of XAI evaluation, the Co-12 properties, and an overview of quantitative evaluation methods \citeref{NAUTA2023}
\item Doshi-Velez and Kim: the three levels of evaluation \citeref{DOSHIVELEZ2017}
\item Molnar et al.: general pitfalls of model-agnostic interpretation methods \citeref{MOLNAR2022}
\item Molnar et al.: PD (and permutation feature importance, later in the course) as estimators of properties of the data-generating process \citeref{MOLNAR2023}
\item Hoffman et al.: metrics for explainable AI from a human-factors perspective \citeref{HOFFMAN2018}
\end{itemizeM}
\end{framev}


{\setbeamertemplate{bibliography item}{}
\begin{frame}[allowframebreaks]{References}
\tiny
\begin{thebibliography}{99}\setlength{\itemsep}{0.25em}
\bibitem{adebayo} \citeref{ADEBAYO2018} J.~Adebayo, J.~Gilmer, M.~Muelly, I.~Goodfellow, M.~Hardt, B.~Kim. Sanity Checks for Saliency Maps. NeurIPS 2018.
\bibitem{alvarez} \citeref{ALVAREZMELIS2018} D.~Alvarez-Melis, T.~S.~Jaakkola. On the Robustness of Interpretability Methods. ICML Workshop on Human Interpretability (WHI) 2018.
\bibitem{apley} \citeref{APLEY2020} D.~W.~Apley, J.~Zhu. Visualizing the Effects of Predictor Variables in Black Box Supervised Learning Models. JRSS-B 82(4), 2020.
\bibitem{baniecki} \citeref{BANIECKI2022} H.~Baniecki, W.~Kretowicz, P.~Biecek. Fooling Partial Dependence via Data Poisoning. ECML PKDD 2022.
\bibitem{bansal} \citeref{BANSAL2021} G.~Bansal, T.~Wu, J.~Zhou, R.~Fok, B.~Nushi, E.~Kamar, M.~T.~Ribeiro, D.~S.~Weld. Does the Whole Exceed its Parts? The Effect of AI Explanations on Complementary Team Performance. CHI 2021.
\bibitem{doshi} \citeref{DOSHIVELEZ2017} F.~Doshi-Velez, B.~Kim. Towards A Rigorous Science of Interpretable Machine Learning. arXiv:1702.08608, 2017.
\bibitem{fisher} \citeref{FISHER2019} A.~Fisher, C.~Rudin, F.~Dominici. All Models are Wrong, but Many are Useful: Learning a Variable's Importance by Studying an Entire Class of Prediction Models Simultaneously. JMLR 20(177), 2019.
\bibitem{freitas} \citeref{FREITAS2014} A.~A.~Freitas. Comprehensible Classification Models: A Position Paper. ACM SIGKDD Explorations 15(1), 2014.
\bibitem{friedman} \citeref{FRIEDMAN2008} J.~H.~Friedman, B.~E.~Popescu. Predictive Learning via Rule Ensembles. Annals of Applied Statistics 2(3), 2008.
\bibitem{goldstein} \citeref{GOLDSTEIN2015} A.~Goldstein, A.~Kapelner, J.~Bleich, E.~Pitkin. Peeking Inside the Black Box: Visualizing Statistical Learning with Plots of Individual Conditional Expectation. JCGS, 2015.
\bibitem{hase} \citeref{HASE2020} P.~Hase, M.~Bansal. Evaluating Explainable AI: Which Algorithmic Explanations Help Users Predict Model Behavior? ACL 2020.
\bibitem{herbinger} \citeref{HERBINGER2024} J.~Herbinger, M.~N.~Wright, T.~Nagler, B.~Bischl, G.~Casalicchio. Decomposing Global Feature Effects Based on Feature Interactions. JMLR 25(381), 2024.
\bibitem{hoffman} \citeref{HOFFMAN2018} R.~R.~Hoffman, S.~T.~Mueller, G.~Klein, J.~Litman. Metrics for Explainable AI: Challenges and Prospects. arXiv:1812.04608, 2018.
\bibitem{jacovi} \citeref{JACOVI2020} A.~Jacovi, Y.~Goldberg. Towards Faithfully Interpretable NLP Systems: How Should We Define and Evaluate Faithfulness? ACL 2020.
\bibitem{kaur} \citeref{KAUR2020} H.~Kaur, H.~Nori, S.~Jenkins, R.~Caruana, H.~Wallach, J.~Wortman Vaughan. Interpreting Interpretability: Understanding Data Scientists' Use of Interpretability Tools for Machine Learning. CHI 2020.
\bibitem{lage} \citeref{LAGE2019} I.~Lage, E.~Chen, J.~He, M.~Narayanan, B.~Kim, S.~J.~Gershman, F.~Doshi-Velez. Human Evaluation of Models Built for Interpretability. AAAI HCOMP 2019.
\bibitem{liu} \citeref{LIU2021} Y.~Liu, S.~Khandagale, C.~White, W.~Neiswanger. Synthetic Benchmarks for Scientific Research in Explainable Machine Learning. arXiv:2106.12543, 2021.
\bibitem{molnarbook} \citeref{MOLNARBOOK} C.~Molnar. Interpretable Machine Learning, Chapter ``Evaluation of Interpretability Methods''. Online book.
\bibitem{molnar2020} \citeref{MOLNAR2020} C.~Molnar, G.~Casalicchio, B.~Bischl. Quantifying Model Complexity via Functional Decomposition for Better Post-hoc Interpretability. ECML PKDD 2019 Workshops, CCIS 1167, 2020.
\bibitem{molnar2022} \citeref{MOLNAR2022} C.~Molnar, G.~K\"onig, J.~Herbinger, T.~Freiesleben, S.~Dandl, C.~A.~Scholbeck, G.~Casalicchio, M.~Grosse-Wentrup, B.~Bischl. General Pitfalls of Model-Agnostic Interpretation Methods for Machine Learning Models. xxAI: Beyond Explainable AI, 2022.
\bibitem{molnar2023} \citeref{MOLNAR2023} C.~Molnar, T.~Freiesleben, G.~K\"onig, G.~Casalicchio, M.~N.~Wright, B.~Bischl. Relating the Partial Dependence Plot and Permutation Feature Importance to the Data Generating Process. xAI 2023, CCIS 1901.
\bibitem{moosbauer} \citeref{MOOSBAUER2021} J.~Moosbauer, J.~Herbinger, G.~Casalicchio, M.~Lindauer, B.~Bischl. Explaining Hyperparameter Optimization via Partial Dependence Plots. NeurIPS 2021.
\bibitem{nadeau} \citeref{NADEAU2003} C.~Nadeau, Y.~Bengio. Inference for the Generalization Error. Machine Learning 52, 2003.
\bibitem{nauta} \citeref{NAUTA2023} M.~Nauta et al. From Anecdotal Evidence to Quantitative Evaluation Methods: A Systematic Review on Evaluating Explainable AI. ACM Computing Surveys, 2023.
\bibitem{poursabzi} \citeref{POURSABZI2021} F.~Poursabzi-Sangdeh, D.~G.~Goldstein, J.~M.~Hofman, J.~Wortman Vaughan, H.~Wallach. Manipulating and Measuring Model Interpretability. CHI 2021.
\bibitem{semenova} \citeref{SEMENOVA2022} L.~Semenova, C.~Rudin, R.~Parr. On the Existence of Simpler Machine Learning Models. FAccT 2022.
\end{thebibliography}
\end{frame}
}



\endlecture
\end{document}